%%%%%%%%%%%%%%%%%%%%%%%%%%%%%%%%%%%%%%%%%%%%%%%%%%%%%%%%%%%%%%%%%%%%%%%%%%%%%%%
% Mathematical Foundations for Kolmogorov-Arnold Networks
% A Continuous Study Record
%
% Author: J Song (Alex)
% Contact: jsong.alex@proton.me
%
% Typeset with KaoNotes (https://github.com/AlexandraSloan3/kaonotes),
% which adapts the kaobook class by Federico Marotta.
%
% Build: xelatex -> bibtex -> xelatex -> xelatex   (or: latexmk main.tex)
%%%%%%%%%%%%%%%%%%%%%%%%%%%%%%%%%%%%%%%%%%%%%%%%%%%%%%%%%%%%%%%%%%%%%%%%%%%%%%%

% !TEX program = xelatex
% !TEX encoding = UTF-8

\documentclass[
	a4paper,
	fontsize=10pt,
	twoside=true,
	open=any,        % avoid inserted blank pages between chapters
	titlepage=false, % compact title, no empty title-back page
	numbers=noenddot,
]{kaobook}

%------------------------------------------------------------------------------
% LANGUAGE
% The body of this book is written entirely in English, so no CJK support is
% loaded. If a Traditional Chinese passage ever becomes unavoidable, enable
% xeCJK here and set \setCJKmainfont{Noto Serif TC} -- do not substitute
% another face.
%------------------------------------------------------------------------------

\usepackage{polyglossia}
\setmainlanguage{english}
\usepackage[english=british]{csquotes}

% Pin the monospaced face so that URLs and identifiers look identical on every
% machine. Libertinus Mono ships with the Libertinus family that kaobook
% already requires, so this adds no new font dependency.
\setmonofont{Libertinus Mono}[Scale=MatchLowercase]

%------------------------------------------------------------------------------
% MATHEMATICS AND REFERENCES
%------------------------------------------------------------------------------

\usepackage[framed=true]{kaotheorems}

% Section-scoped theorem counters need globally unique PDF anchors, one per
% environment kaotheorems declares. Without this, two chapters that each
% happen to contain a "Theorem 1" collide on the same internal PDF anchor.
\providecommand*{\theHtheorem}{}
\renewcommand*{\theHtheorem}{\thesection.\arabic{theorem}}
\providecommand*{\theHproposition}{}
\renewcommand*{\theHproposition}{\thesection.\arabic{proposition}}
\providecommand*{\theHlemma}{}
\renewcommand*{\theHlemma}{\thesection.\arabic{lemma}}
\providecommand*{\theHcorollary}{}
\renewcommand*{\theHcorollary}{\thesection.\arabic{corollary}}
\providecommand*{\theHdefinition}{}
\renewcommand*{\theHdefinition}{\thesection.\arabic{definition}}
\providecommand*{\theHassumption}{}
\renewcommand*{\theHassumption}{\thesection.\arabic{assumption}}
\providecommand*{\theHremark}{}
\renewcommand*{\theHremark}{\thesection.\arabic{remark}}
\providecommand*{\theHexample}{}
\renewcommand*{\theHexample}{\thesection.\arabic{example}}
\providecommand*{\theHexercise}{}
\renewcommand*{\theHexercise}{\thesection.\arabic{exercise}}

\usepackage{kaorefs}

\graphicspath{{figures/}}

%------------------------------------------------------------------------------
% NOTATION SHORTHANDS USED THROUGHOUT THE BOOK
%------------------------------------------------------------------------------

\newcommand{\vect}[1]{\mathbf{#1}}          % column vector
\newcommand{\R}{\mathbb{R}}                 % real numbers
\newcommand{\tp}{^{\mathsf{T}}}             % transpose
\DeclareMathOperator{\Col}{Col}              % column space
\DeclareMathOperator{\Null}{Null}            % nullspace
\newcommand{\E}{\mathbb{E}}                 % expectation
\newcommand{\Var}{\operatorname{Var}}       % variance
\newcommand{\Cov}{\operatorname{Cov}}       % covariance
\newcommand{\Prob}{\operatorname{Pr}}       % probability
% Bold Greek needs the maths font, not the text font used by \mathbf.
\newcommand{\vgreek}[1]{\symbf{#1}}

% Chapter-opening provenance line.
\newcommand{\sources}[1]{%
	\begingroup\small\itshape\noindent Sources: #1\par\endgroup\medskip}

%------------------------------------------------------------------------------
% DOCUMENT METADATA
%------------------------------------------------------------------------------

\begin{document}

\subject{KAN Research Preparation Notes}

\title[Mathematical Foundations for KAN]{Mathematical Foundations for
Kolmogorov--Arnold Networks}
\subtitle{A Continuous Study Record}

\author[J Song (Alex)]{J Song (Alex)}

\date{}

% Set after \title so the full title, not the running short title, reaches the
% PDF metadata.
\hypersetup{
	pdftitle={Mathematical Foundations for Kolmogorov-Arnold Networks},
	pdfauthor={J Song (Alex)},
	pdfsubject={KAN Research Preparation Notes},
	pdfkeywords={linear algebra, calculus, probability, deep learning, KAN},
}

%------------------------------------------------------------------------------

\frontmatter

\maketitle

\begingroup
\setlength{\textheight}{230\hscale}
\etocclasstocstyle
\etocstandardlines
\tableofcontents
\endgroup

%------------------------------------------------------------------------------

% The front matter must not push a blank verso page in front of Part I, and no
% part or chapter opening may insert one either. KOMA forces part pages onto a
% recto by default; these redefinitions turn that off without changing anything
% else about the layout.
\makeatletter
\let\cleardoublepage\clearpage
\@ifundefined{cleardoubleoddpage}{}{\let\cleardoubleoddpage\clearpage}
\@ifundefined{cleardoubleoddstandardpage}{}{\let\cleardoubleoddstandardpage\clearpage}
\@ifundefined{cleardoublestandardpage}{}{\let\cleardoublestandardpage\clearpage}
\makeatother

\mainmatter
\setchapterstyle{kao}

\pagelayout{wide}
\addpart{Linear Algebra}
\pagelayout{margin}

\setchapterpreamble[u]{\margintoc}
\chapter{Vectors, Matrices, and Column Space}
\labch{vectors-and-column-space}

\sources{Strang, \emph{Introduction to Linear Algebra}, 6th ed.\ \cite{strang2023ila},
§1.1--1.4 and §2.1; MIT 18.06 \cite{ocw1806}, Lecture~1.}

Linear algebra begins with a bookkeeping question --- how to store a system of
equations --- and almost immediately turns into a geometric one. This chapter
sets up the notation used for the rest of the book and introduces the two
readings of \(A\vect{x}=\vect{b}\) that everything later depends on.

\section{The matrix form of a linear system}
\labsec{matrix-form}

A system of \(m\) linear equations in \(n\) unknowns can be packed into a single
product,
\[
	A\vect{x}=\vect{b},
\]
where \(A\in\R^{m\times n}\) holds the coefficients, \(\vect{x}\in\R^{n}\) holds
the unknowns, and \(\vect{b}\in\R^{m}\) is the target. Throughout this section I
work with
\[
	A=\begin{bmatrix}
		2 & -1 & 0 \\
		-1 & 2 & -1 \\
		0 & -3 & 4
	\end{bmatrix},
	\qquad
	\vect{b}=\begin{bmatrix}0\\-1\\4\end{bmatrix}.
\]

Read one row at a time, the product unpacks into the three equations
\begin{align*}
	2x-y    &= 0,\\
	-x+2y-z &= -1,\\
	-3y+4z  &= 4.
\end{align*}

\marginnote{Nothing has been solved yet. The matrix form is a change of
notation, not a change of content --- but the notation is what makes the two
pictures below visible.}

Substituting \(\vect{x}=\begin{bmatrix}0&0&1\end{bmatrix}\tp\) satisfies all
three equations at once, so this vector solves the system. The interesting
question is not \emph{whether} this particular vector works, but how one would
have found it, and what the two ways of reading the product suggest.

\section{The row picture: equations as constraints}
\labsec{row-picture}

Each row of a three-variable system describes a plane in \(\R^{3}\). A vector
solves the system exactly when it lies on all three planes simultaneously, so
the solution set is the intersection of the three planes. For the example above
that intersection is the single point \((0,0,1)\).

\begin{remark}
The row picture is a picture of \emph{constraints}. Every additional equation
removes a degree of freedom, and a solution is whatever survives all of the
restrictions at once.
\end{remark}

Its usefulness is bounded by the fact that it has to be drawn. With nine
unknowns the same reasoning would require intersecting hyperplanes in
\(\R^{9}\), which no amount of care will make visible. The algebra is
unaffected --- elimination will still run --- but the geometric intuition stops
paying for itself. That is the first reason to look for a second reading of the
same product.

\section{The column picture: targets as combinations}
\labsec{column-picture}

Write \(\vect{a}_1,\vect{a}_2,\vect{a}_3\) for the columns of \(A\). Expanding
the product by columns instead of by rows gives
\[
	A\vect{x}=x\vect{a}_1+y\vect{a}_2+z\vect{a}_3,
\]
so the system asks for coefficients that mix the three columns into the target:
\[
	x\begin{bmatrix}2\\-1\\0\end{bmatrix}
	+y\begin{bmatrix}-1\\2\\-3\end{bmatrix}
	+z\begin{bmatrix}0\\-1\\4\end{bmatrix}
	=\begin{bmatrix}0\\-1\\4\end{bmatrix}.
\]
In this instance the answer can be read off without any calculation: the target
\emph{is} the third column, so \((x,y,z)=(0,0,1)\). The example is convenient,
but the idea it illustrates is not.

\begin{definition}
A vector \(\vect{b}\) is a \emph{linear combination} of
\(\vect{a}_1,\ldots,\vect{a}_n\) when there exist scalars \(x_1,\ldots,x_n\)
with
\[
	\vect{b}=x_1\vect{a}_1+\cdots+x_n\vect{a}_n .
\]
The set of all such combinations is the \emph{column space} of \(A\), written
\(\Col(A)\).
\end{definition}

The definition converts a question about solving into a question about
membership:
\[
	A\vect{x}=\vect{b}\ \text{is solvable}
	\quad\Longleftrightarrow\quad
	\vect{b}\in\Col(A).
\]

\marginnote{The row picture asks whether a candidate satisfies every
restriction. The column picture asks which recipe of the available ingredients
produces the target. Both are correct; they suggest different algorithms.}

Changing the target tests a different recipe built from the same ingredients.
Taking one of the first column and one of the second,
\[
	\begin{bmatrix}2\\-1\\0\end{bmatrix}
	+\begin{bmatrix}-1\\2\\-3\end{bmatrix}
	=\begin{bmatrix}1\\1\\-3\end{bmatrix},
\]
so the system with target \(\begin{bmatrix}1&1&-3\end{bmatrix}\tp\) is solved by
\(\vect{x}=\begin{bmatrix}1&1&0\end{bmatrix}\tp\).

\subsection{When is every target reachable?}
\labsubsec{reachability}

For a square matrix \(A\in\R^{n\times n}\) the natural question is whether
\(A\vect{x}=\vect{b}\) can be solved for \emph{every} \(\vect{b}\in\R^{n}\). In
the language just introduced this asks whether the columns fill the whole space,
\[
	\Col(A)=\R^{n}.
\]

If the columns are linearly independent they span \(\R^{n}\), every target is
reachable, and the coefficients are unique. The matrix is then invertible and
\(\vect{x}=A^{-1}\vect{b}\).

If instead the columns are dependent, at least one of them contributes no new
direction: it already lies in the span of the others. In \(\R^{3}\) the extreme
case is three columns lying in a common plane. Every combination stays in that
plane, so any target off the plane is unreachable, while targets on it are
reached in infinitely many ways.

\begin{remark}
Dependence and invertibility are the same phenomenon seen from two sides.
Independent columns give exactly one solution for each of the \(n\)-dimensional
family of targets; dependent columns give either none or infinitely many,
depending on where the target sits.
\end{remark}

This dichotomy --- invertible against singular --- is the thread that runs
through elimination in \refch{elimination}, the four subspaces in
\refch{four-subspaces}, and the singular value decomposition in \refch{svd}. It
reappears, transformed, whenever a network layer is asked whether it can
represent a given output.

\section{Matrix--vector multiplication, two readings}
\labsec{matvec}

The column viewpoint supplies the definition of the product that is easiest to
reason with. If
\[
	A=\begin{bmatrix}\vert&&\vert\\ \vect{a}_1&\cdots&\vect{a}_n\\ \vert&&\vert\end{bmatrix},
	\qquad
	\vect{x}=\begin{bmatrix}x_1\\\vdots\\x_n\end{bmatrix},
\]
then
\begin{equation}
	A\vect{x}=x_1\vect{a}_1+\cdots+x_n\vect{a}_n .
	\labeq{matvec-columns}
\end{equation}

The rule taught first in school computes one entry at a time: entry \(i\) of
\(A\vect{x}\) is the dot product of row \(i\) of \(A\) with \(\vect{x}\). Both
give the same numbers. Taking
\[
	A=\begin{bmatrix}2&5\\1&3\end{bmatrix},\qquad
	\vect{x}=\begin{bmatrix}1\\2\end{bmatrix},
\]
the column reading gives
\[
	1\begin{bmatrix}2\\1\end{bmatrix}+2\begin{bmatrix}5\\3\end{bmatrix}
	=\begin{bmatrix}12\\7\end{bmatrix},
\]
and the row reading gives the same vector through \(2\cdot1+5\cdot2=12\) and
\(1\cdot1+3\cdot2=7\).

\begin{remark}
The two readings are not rivals; they answer different questions. Row by row is
how the product is \emph{computed} entry by entry. Column by column is what the
product \emph{means}: \(A\vect{x}\) never leaves the column space of \(A\), no
matter which \(\vect{x}\) is supplied.
\end{remark}

That last observation is worth stating on its own, because it is the reason
\refeq{matvec-columns} is the form used throughout this book.

\begin{proposition}
For every \(\vect{x}\in\R^{n}\) the product \(A\vect{x}\) lies in
\(\Col(A)\), and every element of \(\Col(A)\) arises as \(A\vect{x}\) for some
\(\vect{x}\).
\end{proposition}

\begin{proof}
Equation~\refeq{matvec-columns} exhibits \(A\vect{x}\) as a linear combination
of the columns, which is membership in \(\Col(A)\) by definition. Conversely a
combination \(x_1\vect{a}_1+\cdots+x_n\vect{a}_n\) is the product of \(A\) with
the coefficient vector \(\vect{x}=[x_1,\ldots,x_n]\tp\).
\end{proof}

\section{Span, independence, and the shape of a solution set}
\labsec{span-independence}

Two definitions make the earlier discussion precise.

\begin{definition}
The \emph{span} of \(\vect{a}_1,\ldots,\vect{a}_n\) is the set of all their
linear combinations. The vectors are \emph{linearly independent} when
\[
	x_1\vect{a}_1+\cdots+x_n\vect{a}_n=\vect{0}
	\quad\text{forces}\quad
	x_1=\cdots=x_n=0 ,
\]
and \emph{linearly dependent} otherwise.
\end{definition}

\marginnote{Independence is a statement about the zero target only, yet it
controls the solution count for every target at once. That leverage is what
makes it the central definition of the subject.}

Dependence is exactly the statement that some non-zero \(\vect{x}\) satisfies
\(A\vect{x}=\vect{0}\). Suppose such an \(\vect{x}_0\neq\vect{0}\) exists and
\(\vect{x}_p\) solves \(A\vect{x}=\vect{b}\). Then for every scalar \(c\),
\[
	A(\vect{x}_p+c\vect{x}_0)=A\vect{x}_p+cA\vect{x}_0=\vect{b}+\vect{0}=\vect{b},
\]
so one solution has multiplied into infinitely many. This is the whole
explanation for the trichotomy that a linear system obeys.

\begin{theorem}
A system \(A\vect{x}=\vect{b}\) has no solution, exactly one solution, or
infinitely many. It never has, for instance, exactly two.
\end{theorem}

\begin{proof}
If there is no solution the first case holds. If \(\vect{x}_p\) and
\(\vect{x}_q\) are distinct solutions then \(\vect{x}_0=\vect{x}_p-\vect{x}_q\)
is non-zero and satisfies \(A\vect{x}_0=\vect{0}\); the displayed calculation
above then produces a distinct solution for every scalar \(c\), so there are
infinitely many. Hence two or more solutions already forces infinitely many.
\end{proof}

The set of vectors killed by \(A\) therefore deserves a name of its own; it is
the nullspace, and \refch{four-subspaces} is largely about it.

\section{Notation used in this book}
\labsec{notation}

Lower-case bold letters denote column vectors, capitals denote matrices, and
\(A\tp\) is the transpose. The columns of \(A\) are \(\vect{a}_1,\ldots\); the
entry in row \(i\) and column \(j\) is \(a_{ij}\). The symbol \(\R^{n}\) is the
set of real column vectors of length \(n\), and \(I\) is the identity matrix
whose size is always clear from context.

Vectors are columns unless stated otherwise, so a row vector is written
\(\vect{v}\tp\). Data matrices are the one place where this convention grates:
a batch of \(n\) examples with \(d\) features is stored as \(X\in\R^{n\times d}\)
with one example per \emph{row}, which is the convention every numerical library
uses. Whenever both appear together the shapes are stated explicitly rather than
left to be inferred.

\setchapterpreamble[u]{\margintoc}
\chapter{Solving \texorpdfstring{\(A\vect{x}=\vect{b}\)}{Ax=b}: Elimination and
\texorpdfstring{\(A=LU\)}{A=LU}}
\labch{elimination}

\sources{Strang, \emph{Introduction to Linear Algebra}, 6th ed.\ \cite{strang2023ila},
§1.4 and §2.2--2.7; MIT 18.06 \cite{ocw1806}, Lectures~2--5.}

\refch{vectors-and-column-space} explained what a solution \emph{is}. This
chapter is about how to find one without guessing, and about what the procedure
leaves behind: a factorisation of \(A\) that is worth more than the solution it
was invented to produce.

\section{Elimination on an example}
\labsec{elimination-example}

Take
\[
	A=\begin{bmatrix}1&2&1\\3&8&1\\0&4&1\end{bmatrix},
	\qquad
	\vect{b}=\begin{bmatrix}2\\12\\2\end{bmatrix},
\]
and work with the augmented array \(\begin{bmatrix}A&\vect{b}\end{bmatrix}\) so
that the right-hand side receives every operation the left-hand side does.

The first \emph{pivot} is the entry \(a_{11}=1\). Subtracting \(3\times\)row~1
from row~2 clears the entry below it, and the \((3,1)\) entry is already zero, so
nothing is needed there:
\[
	\left[\begin{array}{ccc|c}
		1&2&1&2\\ 3&8&1&12\\ 0&4&1&2
	\end{array}\right]
	\longrightarrow
	\left[\begin{array}{ccc|c}
		1&2&1&2\\ 0&2&-2&6\\ 0&4&1&2
	\end{array}\right].
\]

\marginnote{The multiplier \(3\) is the number of row~1 subtracted from row~2. It
is worth recording rather than discarding: \refsec{lu} assembles all the
multipliers into a matrix.}

The second pivot is the \(2\) now sitting in position \((2,2)\). Subtracting
twice row~2 from row~3 clears below it:
\[
	\longrightarrow
	\left[\begin{array}{ccc|c}
		1&2&1&2\\ 0&2&-2&6\\ 0&0&5&-10
	\end{array}\right]
	=\left[\begin{array}{c|c}U&\vect{c}\end{array}\right].
\]

The left block \(U\) is \emph{upper triangular}, the pivots are \(1,2,5\), and
the transformed right-hand side is \(\vect{c}=[2,6,-10]\tp\).

\begin{definition}
A \emph{pivot} is the first non-zero entry in a row at the moment that row is
used to clear the entries beneath it. The positions of the pivots, not their
values, determine the structure of the solution set.
\end{definition}

\section{Back-substitution}
\labsec{back-substitution}

Triangular systems are solved from the bottom up. Reading \(U\vect{x}=\vect{c}\)
as equations,
\begin{align*}
	x+2y+z &= 2,\\
	2y-2z  &= 6,\\
	5z     &= -10,
\end{align*}
the last equation gives \(z=-2\). Substituting into the second,
\(2y-2(-2)=6\), so \(2y=2\) and \(y=1\). Substituting both into the first,
\(x+2(1)+(-2)=2\), so \(x=2\). Hence
\[
	\vect{x}=\begin{bmatrix}2\\1\\-2\end{bmatrix},
\]
and substituting into the original \(A\vect{x}\) returns
\([2,\,6+8-2,\,4-2]\tp=[2,12,2]\tp=\vect{b}\), as required.

\marginnote{Always substitute back into the \emph{original} matrix, not into
\(U\). An arithmetic slip made during elimination is invisible to a check
performed on the eliminated system.}

\section{Elimination matrices}
\labsec{elimination-matrices}

Each elimination step is itself a matrix multiplication. Subtracting three
copies of row~1 from row~2 is achieved by left multiplication with
\[
	E_{21}=\begin{bmatrix}1&0&0\\-3&1&0\\0&0&1\end{bmatrix},
\]
which differs from the identity in exactly one off-diagonal entry. Indeed
\(E_{21}A\) reproduces the first array computed above. The second step is
\[
	E_{32}=\begin{bmatrix}1&0&0\\0&1&0\\0&-2&1\end{bmatrix},
	\qquad
	E_{32}\bigl(E_{21}A\bigr)=U .
\]

Because matrix multiplication is associative, the two steps can be combined
before touching \(A\):
\[
	\bigl(E_{32}E_{21}\bigr)A=U .
\]

\begin{remark}
Associativity is doing real work here. It is what allows a sequence of row
operations to be replaced by a single matrix, and it is why the whole
elimination process can be recorded as a factorisation instead of as a script of
instructions.
\end{remark}

Undoing a step is as easy as performing it: if \(E_{21}\) subtracts three copies
of row~1 from row~2, then adding them back inverts it,
\[
	E_{21}^{-1}=\begin{bmatrix}1&0&0\\3&1&0\\0&0&1\end{bmatrix},
	\qquad
	E_{21}^{-1}E_{21}=I .
\]

\section{Matrix multiplication, four ways}
\labsec{four-ways}

Since elimination is multiplication, it is worth collecting the ways a product
\(C=AB\) with \(A\in\R^{m\times n}\) and \(B\in\R^{n\times p}\) can be read. All
four produce the same \(C\); each makes a different fact obvious.

\begin{description}
	\item[Entry by entry] \(c_{ij}=\sum_{k=1}^{n}a_{ik}b_{kj}\), the dot product
	of row \(i\) of \(A\) with column \(j\) of \(B\).
	\item[Column by column] column \(j\) of \(C\) is \(A\vect{b}_j\). Hence
	\emph{every column of \(C\) is a combination of the columns of \(A\)}, and
	\(\Col(AB)\subseteq\Col(A)\).
	\item[Row by row] row \(i\) of \(C\) is \(\vect{a}_i\tp B\), a combination of
	the rows of \(B\). Hence the row space of \(AB\) sits inside the row space of
	\(B\).
	\item[Column times row] \(C=\sum_{k=1}^{n}\vect{a}_k\vect{b}_k\tp\), a sum of
	\(n\) rank-one matrices, where \(\vect{a}_k\) is column \(k\) of \(A\) and
	\(\vect{b}_k\tp\) is row \(k\) of \(B\).
\end{description}

\marginnote{The fourth reading is the one that generalises furthest. The
singular value decomposition in \refch{svd} is precisely a column-times-row
expansion in which the terms are ordered by importance.}

\begin{example}
With
\[
	A=\begin{bmatrix}2&7\\3&8\\4&9\end{bmatrix},
	\qquad
	B=\begin{bmatrix}1&6\\0&0\end{bmatrix},
\]
the column-times-row reading gives
\[
	AB=\begin{bmatrix}2\\3\\4\end{bmatrix}\begin{bmatrix}1&6\end{bmatrix}
	+\begin{bmatrix}7\\8\\9\end{bmatrix}\begin{bmatrix}0&0\end{bmatrix}
	=\begin{bmatrix}2&12\\3&18\\4&24\end{bmatrix}.
\]
The second term vanishes because the second row of \(B\) is zero, which the
entry-by-entry rule would have obscured behind six separate multiplications.
\end{example}

A fifth reading, block multiplication, follows from the same algebra: if \(A\)
and \(B\) are cut into conforming blocks then
\[
	\begin{bmatrix}A_1&A_2\\A_3&A_4\end{bmatrix}
	\begin{bmatrix}B_1&B_2\\B_3&B_4\end{bmatrix}
	=\begin{bmatrix}A_1B_1+A_2B_3 & A_1B_2+A_2B_4\\
	                A_3B_1+A_4B_3 & A_3B_2+A_4B_4\end{bmatrix},
\]
provided the block shapes match. Every practical implementation of matrix
multiplication --- and every batched computation in a neural network --- relies
on this.

\section{Row exchanges and permutation matrices}
\labsec{permutations}

A zero in the pivot position is not fatal on its own. If a row further down has
a non-zero entry in that column, exchanging the two rows restores a usable
pivot. The exchange is again a matrix: for two rows,
\[
	P=\begin{bmatrix}0&1\\1&0\end{bmatrix},
	\qquad
	P\begin{bmatrix}a&b\\c&d\end{bmatrix}=\begin{bmatrix}c&d\\a&b\end{bmatrix}.
\]

\begin{remark}
Multiplying on the \emph{left} rearranges rows; multiplying on the \emph{right}
rearranges columns. The same \(P\) gives
\[
	\begin{bmatrix}a&b\\c&d\end{bmatrix}P=\begin{bmatrix}b&a\\d&c\end{bmatrix}.
\]
Side matters, and the habit of checking which side an operation acts on prevents
a large fraction of avoidable errors.
\end{remark}

\begin{definition}
A \emph{permutation matrix} is obtained from \(I\) by reordering its rows. Every
permutation matrix satisfies \(P^{-1}=P\tp\), because the reordering is undone
by the reverse reordering.
\end{definition}

With exchanges allowed, elimination on an invertible matrix always succeeds, and
the factorisation of \refsec{lu} becomes \(PA=LU\).

\section{When elimination fails}
\labsec{elimination-failure}

Failure occurs only when a pivot position holds zero \emph{and} every entry
below it in that column is zero as well: there is nothing left to exchange with.
A schematic third stage of the form
\[
	\begin{bmatrix}\ast&\ast&\ast\\ 0&\ast&\ast\\ 0&0&0\end{bmatrix}
\]
has only two pivots for three columns. The matrix is then singular, and by
\refsubsec{reachability} some targets are unreachable while others are reached in
infinitely many ways.

\begin{remark}
It is worth separating two situations that look alike. A zero \emph{encountered}
in a pivot position is a bookkeeping accident, repaired by a row exchange. A
zero that \emph{survives} every available exchange is a statement about the
matrix: its columns are dependent, and no algorithm can repair that.
\end{remark}

\section{Inverses and Gauss--Jordan}
\labsec{inverses}

For a square matrix, \(A^{-1}\) is the matrix with \(A^{-1}A=I=AA^{-1}\). It
exists exactly when \(A\) is non-singular. The quickest certificate of
\emph{non}-existence is a non-zero vector in the nullspace.

\begin{proposition}
If \(A\vect{x}=\vect{0}\) for some \(\vect{x}\neq\vect{0}\), then \(A\) has no
inverse.
\end{proposition}

\begin{proof}
Suppose \(A^{-1}\) existed. Multiplying \(A\vect{x}=\vect{0}\) on the left gives
\(\vect{x}=A^{-1}\vect{0}=\vect{0}\), contradicting \(\vect{x}\neq\vect{0}\).
\end{proof}

\begin{example}
For \(A=\begin{bmatrix}1&3\\2&6\end{bmatrix}\) the vector
\(\vect{x}=\begin{bmatrix}3&-1\end{bmatrix}\tp\) gives
\[
	\begin{bmatrix}1&3\\2&6\end{bmatrix}\begin{bmatrix}3\\-1\end{bmatrix}
	=\begin{bmatrix}0\\0\end{bmatrix},
\]
so this \(A\) is singular. The second column is three times the first, which is
the same fact stated in the column picture.
\end{example}

To \emph{find} an inverse, note that column \(j\) of \(A^{-1}\) is the solution
of \(A\vect{x}=\vect{e}_j\), where \(\vect{e}_j\) is column \(j\) of \(I\). The
\(n\) systems share the same left-hand side, so they can be eliminated together.
Gauss--Jordan places \(A\) and \(I\) side by side and reduces until the left
block is the identity:
\[
	\begin{bmatrix}A&I\end{bmatrix}\longrightarrow\begin{bmatrix}I&A^{-1}\end{bmatrix}.
\]

\begin{example}
Take \(A=\begin{bmatrix}1&3\\2&7\end{bmatrix}\). Subtracting \(2\times\)row~1
from row~2, then \(3\times\)the new row~2 from row~1:
\[
	\left[\begin{array}{cc|cc}1&3&1&0\\2&7&0&1\end{array}\right]
	\rightarrow
	\left[\begin{array}{cc|cc}1&3&1&0\\0&1&-2&1\end{array}\right]
	\rightarrow
	\left[\begin{array}{cc|cc}1&0&7&-3\\0&1&-2&1\end{array}\right].
\]
So \(A^{-1}=\begin{bmatrix}7&-3\\-2&1\end{bmatrix}\), and multiplying out
confirms \(AA^{-1}=I\).
\end{example}

\marginnote{Why the method works: the row operations amount to left
multiplication by some matrix \(E\) with \(EA=I\), and that identity says
\(E=A^{-1}\). The right-hand block simply records \(E\) by tracking what the
same operations do to \(I\).}

Two facts about inverses are used constantly. First, the inverse of a product
reverses the order,
\[
	(AB)^{-1}=B^{-1}A^{-1},
\]
which is checked by multiplying: \(AB\,B^{-1}A^{-1}=AIA^{-1}=I\). Second,
inversion commutes with transposition, \((A\tp)^{-1}=(A^{-1})\tp\).

\begin{remark}
In numerical work one almost never forms \(A^{-1}\) explicitly. Solving
\(A\vect{x}=\vect{b}\) by elimination costs about a third as much as inverting
\(A\) and is better behaved in floating point. The inverse is a tool for
reasoning, not usually for computing.
\end{remark}

\section{The factorisation \texorpdfstring{\(A=LU\)}{A=LU}}
\labsec{lu}

Return to \(\bigl(E_{32}E_{21}\bigr)A=U\) and move the elimination matrices to
the other side:
\[
	A=E_{21}^{-1}E_{32}^{-1}\,U=LU .
\]
The remarkable part is what \(L\) contains. Multiplying the two inverses,
\[
	L=\begin{bmatrix}1&0&0\\3&1&0\\0&0&1\end{bmatrix}
	  \begin{bmatrix}1&0&0\\0&1&0\\0&2&1\end{bmatrix}
	 =\begin{bmatrix}1&0&0\\3&1&0\\0&2&1\end{bmatrix},
\]
whose subdiagonal entries are exactly the multipliers \(3\) and \(2\) used
during elimination --- in their original positions, with no interference between
them. Verifying,
\[
	LU=\begin{bmatrix}1&0&0\\3&1&0\\0&2&1\end{bmatrix}
	   \begin{bmatrix}1&2&1\\0&2&-2\\0&0&5\end{bmatrix}
	 =\begin{bmatrix}1&2&1\\3&8&1\\0&4&1\end{bmatrix}=A .
\]

\begin{theorem}
If elimination on \(A\) needs no row exchanges, then \(A=LU\) where \(U\) is the
upper triangular matrix produced by elimination and \(L\) is lower triangular
with ones on the diagonal and the multiplier \(\ell_{ij}\) in position
\((i,j)\). With row exchanges, the same statement holds for a permuted matrix:
\(PA=LU\).
\end{theorem}

\marginnote{The multipliers land in \(L\) unchanged only because the inverses are
multiplied in the reverse order of the eliminations. Multiplying them in the
forward order produces cross terms and destroys the pattern.}

Once \(A=LU\) is available, every subsequent right-hand side is cheap: solve
\(L\vect{c}=\vect{b}\) downwards, then \(U\vect{x}=\vect{c}\) upwards. Both are
triangular, so the expensive part of the work is done once and reused. A
symmetric variant splits the pivots out as \(A=LDU\) with \(D\) diagonal, and for
symmetric positive definite matrices it specialises to \(A=LL\tp\); positive
definiteness is taken up in \refch{eigenvalues}.

\section{Transposes and symmetry}
\labsec{transposes}

The transpose exchanges rows and columns, \((A\tp)_{ij}=a_{ji}\). Its two rules
are
\[
	(AB)\tp=B\tp A\tp,
	\qquad
	(A\tp)\tp=A .
\]
A matrix is \emph{symmetric} when \(A\tp=A\). Symmetric matrices are not a
curiosity: for any \(A\) whatsoever, both \(A\tp A\) and \(AA\tp\) are
symmetric, since
\[
	(A\tp A)\tp=A\tp(A\tp)\tp=A\tp A .
\]

\marginnote{\(A\tp A\) is the matrix behind least squares in
\refch{orthogonality} and behind the singular value decomposition in
\refch{svd}. Almost every time a rectangular matrix has to be turned into a
square one, this is how it is done.}

\section{The cost of elimination}
\labsec{cost}

Clearing the first column of an \(n\times n\) matrix touches roughly \(n^{2}\)
entries, the second about \((n-1)^{2}\), and so on, giving
\[
	n^{2}+(n-1)^{2}+\cdots+1\approx\tfrac{1}{3}n^{3}
\]
operations for the factorisation, while each back-substitution costs about
\(n^{2}\). The cubic term is why reusing \(A=LU\) across many right-hand
sides matters, and why the size of a linear solve grows uncomfortably fast ---
the practical reason that iterative and gradient-based methods dominate once
\(n\) reaches the scale of a machine learning problem.

\setchapterpreamble[u]{\margintoc}
\chapter{The Four Fundamental Subspaces}
\labch{four-subspaces}

\sources{Strang, \emph{Introduction to Linear Algebra}, 6th ed.\ \cite{strang2023ila},
§3.1--3.5; MIT 18.06 \cite{ocw1806}, Lectures~6--11.}

Elimination answers a question about one right-hand side. The subspaces in this
chapter answer the question for all right-hand sides at once, and they do it
with four numbers: two dimensions in the input space and two in the output
space.

\section{Vector spaces and subspaces}
\labsec{subspaces}

\begin{definition}
A \emph{subspace} of \(\R^{n}\) is a non-empty set \(V\subseteq\R^{n}\) closed
under addition and scalar multiplication: if \(\vect{u},\vect{v}\in V\) and
\(c\in\R\) then \(\vect{u}+\vect{v}\in V\) and \(c\vect{v}\in V\).
\end{definition}

Taking \(c=0\) shows that every subspace contains \(\vect{0}\), which is the
quickest test to apply first. A plane through the origin in \(\R^{3}\) is a
subspace; the same plane shifted away from the origin is not, since it fails
that test.

\marginnote{Closure is the whole content of the definition. It says that the two
operations the subject is built from cannot carry you out of the set --- which
is exactly what makes linear questions answerable inside it.}

The subspaces that matter here are not chosen arbitrarily. Each of the four is
attached to a matrix and describes one aspect of what that matrix does.

\section{The column space}
\labsec{column-space}

The column space \(\Col(A)\subseteq\R^{m}\) was defined in
\refch{vectors-and-column-space} as the set of all combinations of the columns.
It is a subspace because a combination of combinations is again a combination,
and it settles solvability: \(A\vect{x}=\vect{b}\) has a solution precisely when
\(\vect{b}\in\Col(A)\).

Throughout this chapter I use
\[
	A=\begin{bmatrix}
		1&2&2&2\\
		2&4&6&8\\
		3&6&8&10
	\end{bmatrix}\in\R^{3\times4}.
\]
Its second column is twice its first, so the columns are dependent from the
start and the column space is smaller than \(\R^{3}\).

\section{The nullspace}
\labsec{nullspace}

\begin{definition}
The \emph{nullspace} of \(A\in\R^{m\times n}\) is
\[
	\Null(A)=\{\vect{x}\in\R^{n}: A\vect{x}=\vect{0}\}\subseteq\R^{n}.
\]
\end{definition}

It is a subspace: if \(A\vect{u}=\vect{0}\) and \(A\vect{v}=\vect{0}\) then
\(A(\vect{u}+c\vect{v})=\vect{0}\). Unlike the column space it lives in the
\emph{input} space, and it is exactly the object that \refsec{span-independence}
showed controls how many solutions a system has.

To compute it, eliminate. Subtracting \(2\times\)row~1 from row~2 and
\(3\times\)row~1 from row~3, then subtracting the new row~2 from the new row~3,
\[
	A\longrightarrow
	\begin{bmatrix}1&2&2&2\\0&0&2&4\\0&0&2&4\end{bmatrix}
	\longrightarrow
	U=\begin{bmatrix}1&2&2&2\\0&0&2&4\\0&0&0&0\end{bmatrix}.
\]

Two pivots appear, in columns~1 and~3. The row of zeros is not an accident: it
records that row~3 of \(A\) was the sum of rows~1 and~2 all along.

\begin{definition}
Columns containing a pivot are \emph{pivot columns}; the rest are \emph{free
columns}. The unknowns attached to free columns are \emph{free variables},
because they may be assigned any values whatever, after which the pivot
variables are determined.
\end{definition}

Here \(x_2\) and \(x_4\) are free. Setting \(x_2=1,x_4=0\) and back-substituting
gives \(x_3=0\) and \(x_1=-2\); setting \(x_2=0,x_4=1\) gives \(x_3=-2\) and
\(x_1=2\). So
\[
	\vect{s}_1=\begin{bmatrix}-2\\1\\0\\0\end{bmatrix},
	\qquad
	\vect{s}_2=\begin{bmatrix}2\\0\\-2\\1\end{bmatrix},
	\qquad
	\Null(A)=\operatorname{span}\{\vect{s}_1,\vect{s}_2\}.
\]

\marginnote{Check both by hand. \(A\vect{s}_1=-2\vect{a}_1+\vect{a}_2\) and the
second column is twice the first, so the result is \(\vect{0}\). Similarly
\(A\vect{s}_2=2\vect{a}_1-2\vect{a}_3+\vect{a}_4=\vect{0}\).}

\section{Rank, echelon form, and \texorpdfstring{\(A=CR\)}{A=CR}}
\labsec{rank-cr}

\begin{definition}
The \emph{rank} \(r\) of \(A\) is the number of pivots produced by elimination.
\end{definition}

For the running example \(r=2\). Continuing elimination until each pivot is~1
and the entries above it are cleared produces the \emph{reduced row echelon
form}
\[
	R_0=\begin{bmatrix}1&2&0&-2\\0&0&1&2\\0&0&0&0\end{bmatrix}.
\]

Discarding the zero row leaves \(R\in\R^{2\times4}\), and collecting the pivot
columns of the original matrix gives
\[
	C=\begin{bmatrix}1&2\\2&6\\3&8\end{bmatrix},
	\qquad
	R=\begin{bmatrix}1&2&0&-2\\0&0&1&2\end{bmatrix}.
\]

\begin{theorem}
Every matrix factors as \(A=CR\), where the \(r\) columns of \(C\) are the pivot
columns of \(A\) and the \(r\) rows of \(R\) are the non-zero rows of the reduced
row echelon form.
\end{theorem}

Multiplying out confirms it here: the first column of \(CR\) is
\(1\cdot\vect{c}_1=[1,2,3]\tp\), the second is \(2\vect{c}_1=[2,4,6]\tp\), the
third is \(\vect{c}_2=[2,6,8]\tp\), and the fourth is
\(-2\vect{c}_1+2\vect{c}_2=[2,8,10]\tp\).

\begin{remark}
\(A=CR\) is the cleanest statement of what rank means. The columns of \(A\) are
combinations of \(r\) independent columns, and the rows of \(A\) are
combinations of \(r\) independent rows --- the \emph{same} \(r\) in both cases.
Row rank and column rank being equal is often presented as a surprise; here it
is a consequence of a single factorisation.
\end{remark}

\section{The complete solution of \texorpdfstring{\(A\vect{x}=\vect{b}\)}{Ax=b}}
\labsec{complete-solution}

Take \(\vect{b}=[1,6,7]\tp\). Solving with the free variables set to zero gives a
\emph{particular} solution: the pivot variables must satisfy
\(x_1\vect{a}_1+x_3\vect{a}_3=\vect{b}\), that is \(x_1+2x_3=1\) and
\(2x_1+6x_3=6\), whence \(x_3=2\) and \(x_1=-3\); the third equation
\(3x_1+8x_3=-9+16=7\) is consistent. So
\[
	\vect{x}_p=\begin{bmatrix}-3\\0\\2\\0\end{bmatrix},
	\qquad
	\vect{x}=\vect{x}_p+c_1\vect{s}_1+c_2\vect{s}_2\ \ \text{for all }c_1,c_2\in\R.
\]

\begin{proposition}
If \(A\vect{x}=\vect{b}\) is solvable, its solution set is
\(\vect{x}_p+\Null(A)\): one particular solution shifted by the whole nullspace.
\end{proposition}

\begin{proof}
Any solution \(\vect{x}\) satisfies \(A(\vect{x}-\vect{x}_p)=\vect{b}-\vect{b}
=\vect{0}\), so \(\vect{x}-\vect{x}_p\in\Null(A)\). Conversely adding any
nullspace vector to \(\vect{x}_p\) leaves the product unchanged.
\end{proof}

\marginnote{This is the same statement as the trichotomy in
\refch{vectors-and-column-space}, now with the geometry supplied: the solution
set is a shifted copy of a subspace, so it is empty, a single point, or an
unbounded flat piece.}

\section{Independence, basis, and dimension}
\labsec{basis-dimension}

\begin{definition}
A \emph{basis} of a subspace \(V\) is a set of vectors that is linearly
independent and spans \(V\). The number of vectors in a basis is the
\emph{dimension} of \(V\), written \(\dim V\).
\end{definition}

That the number is the same for every basis is the fact that makes dimension
well defined. For the running example, \(\{\vect{s}_1,\vect{s}_2\}\) is a basis
of the nullspace and the two pivot columns form a basis of the column space, so
\[
	\dim\Col(A)=2,
	\qquad
	\dim\Null(A)=2 .
\]

\section{The four subspaces}
\labsec{four}

Applying the two constructions to \(A\) and to \(A\tp\) yields four subspaces in
total.

\begin{description}
	\item[Column space \(\Col(A)\subseteq\R^{m}\)] spanned by the columns;
	dimension \(r\).
	\item[Nullspace \(\Null(A)\subseteq\R^{n}\)] vectors killed by \(A\);
	dimension \(n-r\).
	\item[Row space \(\Col(A\tp)\subseteq\R^{n}\)] spanned by the rows;
	dimension \(r\).
	\item[Left nullspace \(\Null(A\tp)\subseteq\R^{m}\)] the vectors
	\(\vect{y}\) satisfying \(A\tp\vect{y}=\vect{0}\); dimension \(m-r\).
\end{description}

For the running example \(m=3\), \(n=4\), \(r=2\), so the dimensions are
\(2,2,2,1\). The left nullspace is spanned by \([1,1,-1]\tp\), which encodes the
dependency noticed during elimination: row~1 plus row~2 equals row~3.

\begin{theorem}[Rank--nullity]
For \(A\in\R^{m\times n}\) of rank \(r\),
\[
	\dim\Col(A)+\dim\Null(A)=r+(n-r)=n .
\]
\end{theorem}

\begin{proof}
Each of the \(n\) columns is either a pivot column or a free column. The pivot
columns form a basis of \(\Col(A)\), giving \(r\); each free column contributes
exactly one special solution, and those are independent because each carries a
\(1\) in a position where the others carry \(0\), giving \(n-r\).
\end{proof}

\begin{remark}
The counting has a plain reading. The input space \(\R^{n}\) splits its \(n\)
degrees of freedom into \(r\) that the matrix records and \(n-r\) that it
discards. Nothing else can happen, and no amount of cleverness in choosing
\(\vect{b}\) will change the split.
\end{remark}

\refch{orthogonality} sharpens this picture considerably: the two subspaces
inside \(\R^{n}\) are not merely complementary in dimension, they are orthogonal
complements, and likewise for the two inside \(\R^{m}\).

\setchapterpreamble[u]{\margintoc}
\chapter{Orthogonality, Projection, and Least Squares}
\labch{orthogonality}

\sources{Strang, \emph{Introduction to Linear Algebra}, 6th ed.\ \cite{strang2023ila},
§4.1--4.4; MIT 18.06 \cite{ocw1806}, Lectures~14--17.}

Most systems that arise from measurements have no solution: there are more
equations than unknowns and the data are noisy. Rather than give up, one asks
for the vector that comes closest. Orthogonality is what makes ``closest''
computable.

\section{Inner products, length, and angle}
\labsec{inner-products}

For \(\vect{u},\vect{v}\in\R^{n}\) the inner product is
\(\vect{u}\tp\vect{v}=\sum_i u_iv_i\), the length is
\(\lVert\vect{v}\rVert=\sqrt{\vect{v}\tp\vect{v}}\), and the two vectors are
\emph{orthogonal} when \(\vect{u}\tp\vect{v}=0\).

\begin{proposition}
\(\vect{u}\tp\vect{v}=0\) if and only if
\(\lVert\vect{u}+\vect{v}\rVert^{2}=\lVert\vect{u}\rVert^{2}+\lVert\vect{v}\rVert^{2}\).
\end{proposition}

\begin{proof}
Expanding,
\(\lVert\vect{u}+\vect{v}\rVert^{2}
=(\vect{u}+\vect{v})\tp(\vect{u}+\vect{v})
=\lVert\vect{u}\rVert^{2}+2\vect{u}\tp\vect{v}+\lVert\vect{v}\rVert^{2}\).
The two sides agree exactly when the cross term vanishes.
\end{proof}

\marginnote{So Pythagoras is not an analogy imported from geometry; it is the
algebraic identity above, and it holds in \(\R^{n}\) for every \(n\).}

\section{Orthogonal subspaces}
\labsec{orthogonal-subspaces}

Two subspaces \(V,W\) are orthogonal when \(\vect{v}\tp\vect{w}=0\) for every
\(\vect{v}\in V\) and \(\vect{w}\in W\). The four subspaces of
\refch{four-subspaces} pair up in exactly this way.

\begin{theorem}
For any \(A\in\R^{m\times n}\), the nullspace and the row space are orthogonal
complements in \(\R^{n}\), and the left nullspace and the column space are
orthogonal complements in \(\R^{m}\).
\end{theorem}

\begin{proof}
If \(A\vect{x}=\vect{0}\) then every row of \(A\) has zero inner product with
\(\vect{x}\), hence so does every combination of rows: \(\Null(A)\) is orthogonal
to \(\Col(A\tp)\). Their dimensions are \(n-r\) and \(r\), which sum to \(n\), so
neither can be enlarged inside the other's complement. Applying the same
argument to \(A\tp\) gives the second pair.
\end{proof}

\begin{remark}
This upgrades rank--nullity from an accounting identity to a geometric
statement. The input space splits into a part the matrix reproduces faithfully
and a part it annihilates, and the two parts meet at right angles.
\end{remark}

\section{Projection onto a line}
\labsec{projection-line}

Given \(\vect{a}\neq\vect{0}\), the point of the line \(\{c\vect{a}\}\) closest
to \(\vect{b}\) is the one for which the error \(\vect{b}-c\vect{a}\) is
perpendicular to \(\vect{a}\):
\[
	\vect{a}\tp(\vect{b}-c\vect{a})=0
	\quad\Longrightarrow\quad
	c=\frac{\vect{a}\tp\vect{b}}{\vect{a}\tp\vect{a}},
	\qquad
	\vect{p}=\frac{\vect{a}\tp\vect{b}}{\vect{a}\tp\vect{a}}\,\vect{a}.
\]

Writing this as a matrix acting on \(\vect{b}\) gives the projection matrix
\[
	P=\frac{\vect{a}\vect{a}\tp}{\vect{a}\tp\vect{a}},
	\qquad
	P^{2}=P,\qquad P\tp=P .
\]

\marginnote{\(P^{2}=P\) says projecting twice is projecting once, which any
sensible notion of ``closest point'' must satisfy. Symmetry is the algebraic
trace of the right angle.}

\begin{example}
With \(\vect{a}=[1,1,1]\tp\) and \(\vect{b}=[1,2,2]\tp\),
\(\vect{a}\tp\vect{b}=5\) and \(\vect{a}\tp\vect{a}=3\), so
\(\vect{p}=\tfrac{5}{3}[1,1,1]\tp\) and the error
\(\vect{b}-\vect{p}=\tfrac{1}{3}[-2,1,1]\tp\) indeed satisfies
\(\vect{a}\tp(\vect{b}-\vect{p})=0\).
\end{example}

\section{Projection onto a subspace}
\labsec{projection-subspace}

Replace the single vector by the columns of \(A\in\R^{m\times n}\) and look for
\(\hat{\vect{x}}\) making \(\vect{b}-A\hat{\vect{x}}\) orthogonal to every
column. That is \(A\tp(\vect{b}-A\hat{\vect{x}})=\vect{0}\), or

\begin{equation}
	A\tp A\hat{\vect{x}}=A\tp\vect{b},
	\labeq{normal-equations}
\end{equation}

the \emph{normal equations}. When the columns of \(A\) are independent,
\(A\tp A\) is invertible and
\[
	\hat{\vect{x}}=(A\tp A)^{-1}A\tp\vect{b},
	\qquad
	\vect{p}=A\hat{\vect{x}}=A(A\tp A)^{-1}A\tp\vect{b}.
\]

\begin{proposition}
If the columns of \(A\) are linearly independent then \(A\tp A\) is invertible.
\end{proposition}

\begin{proof}
Suppose \(A\tp A\vect{x}=\vect{0}\). Then
\(\vect{x}\tp A\tp A\vect{x}=\lVert A\vect{x}\rVert^{2}=0\), so
\(A\vect{x}=\vect{0}\), and independence of the columns forces
\(\vect{x}=\vect{0}\). A square matrix with trivial nullspace is invertible.
\end{proof}

\marginnote{The proof is the reason \(A\tp A\) appears everywhere. It converts a
rectangular problem into a square one without losing information, provided the
columns were independent to begin with.}

\section{Least squares}
\labsec{least-squares}

Solving \refeq{normal-equations} minimises \(\lVert A\vect{x}-\vect{b}\rVert^{2}\)
over all \(\vect{x}\). This is the standard way to fit a model with more
observations than parameters.

\begin{example}
Fit a straight line \(b=C+Dt\) through the points \((1,1)\), \((2,2)\),
\((3,2)\). Then
\[
	A=\begin{bmatrix}1&1\\1&2\\1&3\end{bmatrix},
	\qquad
	\vect{b}=\begin{bmatrix}1\\2\\2\end{bmatrix},
	\qquad
	A\tp A=\begin{bmatrix}3&6\\6&14\end{bmatrix},
	\qquad
	A\tp\vect{b}=\begin{bmatrix}5\\11\end{bmatrix}.
\]
The normal equations \(3C+6D=5\) and \(6C+14D=11\) give \(D=\tfrac{1}{2}\) and
\(C=\tfrac{2}{3}\), so the best line is \(b=\tfrac{2}{3}+\tfrac{1}{2}t\). The
residuals are \(-\tfrac{1}{6},\tfrac{1}{3},-\tfrac{1}{6}\); they sum to zero and
satisfy \(A\tp\vect{e}=\vect{0}\), confirming that the error is orthogonal to the
column space.
\end{example}

\begin{remark}
Least squares is the first place where a linear algebra identity and a
minimisation problem coincide. The same coincidence returns in
\refch{optimization} for quadratic objectives and again in \refch{estimation},
where minimising squared error turns out to be maximum likelihood under Gaussian
noise.
\end{remark}

\section{Orthonormal bases and Gram--Schmidt}
\labsec{gram-schmidt}

A set \(\{\vect{q}_1,\ldots,\vect{q}_n\}\) is \emph{orthonormal} when
\(\vect{q}_i\tp\vect{q}_j\) equals \(1\) for \(i=j\) and \(0\) otherwise.
Collecting them as columns of \(Q\) gives \(Q\tp Q=I\), and when \(Q\) is square
this is \(Q^{-1}=Q\tp\).

Orthonormal columns simplify everything above: the normal equations collapse to
\(\hat{\vect{x}}=Q\tp\vect{b}\), and the projection matrix becomes
\(P=QQ\tp\), with no inverse to compute.

\marginnote{Numerically this is the difference between a stable computation and
an unstable one. Forming \(A\tp A\) squares the condition number of \(A\);
orthogonalising first avoids that.}

Gram--Schmidt manufactures an orthonormal basis from any independent set by
subtracting, from each vector in turn, its projection onto everything already
accepted.

\begin{example}
For \(A=\begin{bmatrix}1&1\\1&0\\0&1\end{bmatrix}\), start with
\(\vect{q}_1=\tfrac{1}{\sqrt2}[1,1,0]\tp\). Subtracting the projection of the
second column,
\[
	\begin{bmatrix}1\\0\\1\end{bmatrix}
	-\tfrac{1}{2}\begin{bmatrix}1\\1\\0\end{bmatrix}
	=\begin{bmatrix}\tfrac12\\-\tfrac12\\1\end{bmatrix},
	\qquad
	\vect{q}_2=\tfrac{1}{\sqrt6}\begin{bmatrix}1\\-1\\2\end{bmatrix}.
\]
\end{example}

\section{The factorisation \texorpdfstring{\(A=QR\)}{A=QR}}
\labsec{qr}

Recording the coefficients used during Gram--Schmidt gives a second
factorisation alongside \(A=LU\).

\begin{theorem}
If \(A\in\R^{m\times n}\) has independent columns then \(A=QR\), where the
columns of \(Q\in\R^{m\times n}\) are orthonormal and \(R\in\R^{n\times n}\) is
upper triangular with positive diagonal entries.
\end{theorem}

For the example above,
\[
	R=\begin{bmatrix}\sqrt2 & \tfrac{1}{\sqrt2}\\[2pt] 0 & \sqrt{\tfrac32}\end{bmatrix},
\]
and multiplying \(QR\) returns the two original columns. Substituting \(A=QR\)
into the normal equations turns them into the triangular system
\(R\hat{\vect{x}}=Q\tp\vect{b}\), which is how least squares is solved in
practice.

\begin{remark}
Two factorisations, two purposes. \(A=LU\) records elimination and answers
questions about solving; \(A=QR\) records orthogonalisation and answers
questions about approximating. \refch{svd} produces a third that does both and
survives rank deficiency as well.
\end{remark}

\setchapterpreamble[u]{\margintoc}
\chapter{Determinants}
\labch{determinants}

\sources{Strang, \emph{Introduction to Linear Algebra}, 6th ed.\ \cite{strang2023ila},
§5.1--5.3; MIT 18.06 \cite{ocw1806}, Lectures~18--20.}

The determinant compresses a square matrix into one number that answers a
yes-or-no question: is the matrix invertible? Its main use in this book is as
the machinery behind eigenvalues in \refch{eigenvalues}.

\section{Three properties that define it}
\labsec{det-properties}

Rather than start from a formula, take three requirements and derive everything
from them. For \(A\in\R^{n\times n}\), the determinant \(\det A\) is the unique
function satisfying:

\begin{enumerate}
	\item \(\det I=1\);
	\item exchanging two rows reverses the sign;
	\item it is linear in each row separately, the other rows being held fixed.
\end{enumerate}

\marginnote{Property~3 is linearity \emph{one row at a time}, not linearity in
the matrix. Multiplying the whole of an \(n\times n\) matrix by \(c\) multiplies
the determinant by \(c^{n}\), once per row.}

\section{What follows immediately}
\labsec{det-consequences}

\begin{proposition}
If two rows of \(A\) are equal then \(\det A=0\).
\end{proposition}

\begin{proof}
Exchanging the two equal rows leaves \(A\) unchanged but reverses the sign by
property~2, so \(\det A=-\det A\).
\end{proof}

\begin{proposition}
Subtracting a multiple of one row from another leaves the determinant unchanged.
\end{proposition}

\begin{proof}
By linearity in the affected row, the determinant of the result is
\(\det A\) minus \(c\) times the determinant of a matrix with two equal rows,
and that second term is zero.
\end{proof}

These two facts make elimination a legitimate way to compute a determinant: the
subtraction steps cost nothing and only the row exchanges have to be counted.

\begin{proposition}
A triangular matrix has determinant equal to the product of its diagonal
entries; a matrix with a zero row has determinant \(0\).
\end{proposition}

\section{Determinant from the pivots}
\labsec{det-pivots}

Combining the last two sections: eliminate \(A\) into \(U\), keeping track of how
many row exchanges were used. Then
\[
	\det A=(-1)^{\,\text{number of exchanges}}\prod_{i=1}^{n}u_{ii},
\]
the product of the pivots up to sign.

\begin{example}
The elimination carried out in \refsec{elimination-example} on
\[
	A=\begin{bmatrix}1&2&1\\3&8&1\\0&4&1\end{bmatrix}
\]
needed no row exchanges and produced the pivots \(1,2,5\), so \(\det A=10\).
\end{example}

\begin{theorem}
\(A\) is invertible if and only if \(\det A\neq0\).
\end{theorem}

\begin{proof}
Elimination produces \(n\) non-zero pivots exactly when \(A\) is non-singular,
and the displayed formula makes \(\det A\) non-zero precisely in that case.
\end{proof}

The multiplicative rule \(\det(AB)=(\det A)(\det B)\) follows from the same
circle of ideas, and gives \(\det(A^{-1})=1/\det A\) as a special case. Transposition also leaves it alone,
\(\det(A\tp)=\det A\), so each statement about rows has a counterpart about
columns.

\section{Cofactors and the big formula}
\labsec{cofactors}

Expanding property~3 across all rows at once gives a sum over the \(n!\)
permutations, one term for each way of picking a single entry from every row and
column. For \(n=3\) this is the familiar six-term expression; for \(n=10\) it has
over three million terms and is useless for computation.

The practical repackaging is cofactor expansion along a row,
\[
	\det A=\sum_{j=1}^{n}a_{ij}C_{ij},
	\qquad
	C_{ij}=(-1)^{i+j}\det M_{ij},
\]
where \(M_{ij}\) deletes row \(i\) and column \(j\).

\begin{example}
Expanding the same matrix along its first row,
\[
	\det A=1(8-4)-2(3-0)+1(12-0)=4-6+12=10,
\]
agreeing with the pivot calculation.
\end{example}

\marginnote{Two routes, one answer. Cofactors are convenient for small or sparse
matrices and for symbolic work; pivots are what any implementation actually
uses.}

\section{Three uses}
\labsec{det-uses}

\paragraph{The inverse formula.}
\(A^{-1}=\dfrac{1}{\det A}\,C\tp\), where \(C\) is the matrix of cofactors. It
makes the dependence of \(A^{-1}\) on the entries of \(A\) explicit, and shows
why a determinant near zero makes the inverse enormous.

\paragraph{Cramer's rule.}
The \(j\)th component of the solution of \(A\vect{x}=\vect{b}\) is
\(x_j=\det B_j/\det A\), where \(B_j\) is \(A\) with column \(j\) replaced by
\(\vect{b}\). Elegant, and far too slow to use.

\paragraph{Volume.}
\(|\det A|\) is the volume of the parallelepiped spanned by the rows of \(A\).
A zero determinant means the solid is flat --- the rows lie in a lower
dimensional subspace --- which is the geometric form of singularity.

\begin{remark}
The volume reading is the one that carries into calculus. In
\refch{jacobian-hessian} the determinant of the Jacobian is the local
volume-scaling factor of a nonlinear map, and it is what appears when a change of
variables is made inside an integral.
\end{remark}

\setchapterpreamble[u]{\margintoc}
\chapter{Eigenvalues and Eigenvectors}
\labch{eigenvalues}

\sources{Strang, \emph{Introduction to Linear Algebra}, 6th ed.\ \cite{strang2023ila},
§6.1--6.5; MIT 18.06 \cite{ocw1806}, Lectures~21--22 and 26--28.}

Everything so far treated \(A\) as a machine for solving. This chapter asks a
different question: are there directions in which \(A\) does nothing more
interesting than stretching? The answer reorganises the matrix around those
directions and makes repeated application easy to describe.

\section{The defining equation}
\labsec{eigen-definition}

\begin{definition}
A non-zero vector \(\vect{x}\) is an \emph{eigenvector} of \(A\in\R^{n\times n}\)
with \emph{eigenvalue} \(\lambda\) when
\[
	A\vect{x}=\lambda\vect{x}.
\]
\end{definition}

Rewriting as \((A-\lambda I)\vect{x}=\vect{0}\) shows that \(\lambda\) is an
eigenvalue exactly when \(A-\lambda I\) is singular, since a non-zero vector must
sit in its nullspace. By \refch{determinants} that is the condition
\[
	\det(A-\lambda I)=0,
\]
the \emph{characteristic equation}. Its left-hand side is a polynomial of degree
\(n\) in \(\lambda\).

\marginnote{Note the order of work: first find the \(\lambda\) that make the
matrix singular, then find the nullspace for each. Eigenvectors are never found
first.}

\begin{example}
For \(A=\begin{bmatrix}2&1\\1&2\end{bmatrix}\),
\[
	\det\begin{bmatrix}2-\lambda&1\\1&2-\lambda\end{bmatrix}
	=(2-\lambda)^{2}-1=\lambda^{2}-4\lambda+3,
\]
with roots \(\lambda_1=3\) and \(\lambda_2=1\). Solving
\((A-3I)\vect{x}=\vect{0}\) gives \(\vect{x}_1=[1,1]\tp\), and
\((A-I)\vect{x}=\vect{0}\) gives \(\vect{x}_2=[1,-1]\tp\). The two eigenvectors
are orthogonal, which \refsec{spectral} explains.
\end{example}

\section{Trace and determinant}
\labsec{trace-det}

Two shortcuts follow from comparing coefficients in the characteristic
polynomial.

\begin{proposition}
The eigenvalues of \(A\), counted with multiplicity, satisfy
\[
	\sum_{i}\lambda_i=\operatorname{tr}A=\sum_i a_{ii},
	\qquad
	\prod_i\lambda_i=\det A .
\]
\end{proposition}

For the example, \(3+1=4=2+2\) and \(3\cdot1=3=4-1\). In two dimensions these two
equations determine the eigenvalues outright, which is often faster than forming
the polynomial.

\begin{remark}
Singularity and a zero eigenvalue are the same statement, since the product of
the eigenvalues is the determinant. The nullspace of \(A\) is precisely the
eigenspace for \(\lambda=0\).
\end{remark}

\section{Diagonalisation and powers}
\labsec{diagonalisation}

Suppose \(A\) has \(n\) independent eigenvectors. Placing them in the columns of
\(X\) and the eigenvalues on the diagonal of \(\Lambda\),
\[
	AX=X\Lambda
	\quad\Longrightarrow\quad
	A=X\Lambda X^{-1}.
\]

The payoff is that powers telescope:
\[
	A^{k}=X\Lambda^{k}X^{-1},
\]
because the inner \(X^{-1}X\) factors cancel. Raising a diagonal matrix to a
power is done entry by entry, so the long-run behaviour of \(A^{k}\) is decided
entirely by the sizes \(|\lambda_i|\).

\begin{example}
Let \(A=\begin{bmatrix}0.8&0.3\\0.2&0.7\end{bmatrix}\), whose columns sum to~1.
Its trace is \(1.5\) and its determinant is \(0.5\), so the eigenvalues solve
\(\lambda^{2}-1.5\lambda+0.5=0\), giving \(\lambda_1=1\) and \(\lambda_2=0.5\).
The eigenvector for \(\lambda_1=1\) is \([3,2]\tp\); normalising it to sum to~1
gives \([0.6,0.4]\tp\). Since \(0.5^{k}\to0\), every starting distribution
converges to that vector.
\end{example}

\marginnote{This is the entire theory of Markov chains in miniature: an
eigenvalue of exactly~1 supplies the steady state, and the second largest
modulus sets the rate of approach.}

Diagonalisation can fail. The matrix \(\begin{bmatrix}0&1\\0&0\end{bmatrix}\)
has \(\lambda=0\) twice but only a one-dimensional eigenspace, so no basis of
eigenvectors exists. Distinct eigenvalues guarantee independence and therefore
diagonalisability; repeated eigenvalues may or may not.

\section{Symmetric matrices}
\labsec{spectral}

Symmetric matrices are the well-behaved case, and they are the case that occurs
in practice --- covariance matrices, Gram matrices \(A\tp A\), and Hessians are
all symmetric.

\begin{theorem}[Spectral theorem]
If \(A\tp=A\) then all eigenvalues of \(A\) are real, eigenvectors for distinct
eigenvalues are orthogonal, and \(A\) can be written
\[
	A=Q\Lambda Q\tp
\]
with \(Q\) orthogonal and \(\Lambda\) real diagonal.
\end{theorem}

\begin{proof}[Orthogonality]
Let \(A\vect{x}=\lambda\vect{x}\) and \(A\vect{y}=\mu\vect{y}\) with
\(\lambda\neq\mu\). Then
\[
	\lambda\,\vect{y}\tp\vect{x}=\vect{y}\tp A\vect{x}
	=(A\vect{y})\tp\vect{x}=\mu\,\vect{y}\tp\vect{x},
\]
using symmetry in the middle step. Hence
\((\lambda-\mu)\vect{y}\tp\vect{x}=0\), and since \(\lambda\neq\mu\) the inner
product vanishes.
\end{proof}

Because \(Q^{-1}=Q\tp\), the factorisation writes \(A\) as a sum of orthogonal
rank-one pieces,
\[
	A=\sum_{i=1}^{n}\lambda_i\,\vect{q}_i\vect{q}_i\tp ,
\]
which is the column-times-row reading of \refsec{four-ways}, applied here to
the basis of eigenvectors.

\section{Positive definite matrices}
\labsec{positive-definite}

\begin{definition}
A symmetric \(A\) is \emph{positive definite} when \(\vect{x}\tp A\vect{x}>0\)
for every \(\vect{x}\neq\vect{0}\), and \emph{positive semidefinite} when the
inequality is weak.
\end{definition}

\begin{proposition}
For a symmetric matrix \(A\) the following are equivalent: every eigenvalue is
positive; every pivot is positive; \(\vect{x}\tp A\vect{x}>0\) for every
\(\vect{x}\neq\vect{0}\); and \(A=B\tp B\) for some \(B\) with independent
columns.
\end{proposition}

\begin{example}
\(A=\begin{bmatrix}2&1\\1&2\end{bmatrix}\) has eigenvalues \(3\) and \(1\),
pivots \(2\) and \(\tfrac32\), and
\(\vect{x}\tp A\vect{x}=2x_1^{2}+2x_1x_2+2x_2^{2}=(x_1+x_2)^{2}+x_1^{2}+x_2^{2}\),
positive unless both components vanish. All three tests agree.
\end{example}

\marginnote{Positive definiteness is the matrix version of ``strictly convex''.
\refch{gradient-methods} uses exactly this: a quadratic with a positive definite
Hessian has a single minimum, and the ratio of largest to smallest eigenvalue
predicts how slowly gradient descent will find it.}

\section{Similar matrices}
\labsec{similar}

Matrices \(A\) and \(M^{-1}AM\) are called \emph{similar}. They have the same
characteristic polynomial and hence the same eigenvalues, because
\[
	\det(M^{-1}AM-\lambda I)=\det\bigl(M^{-1}(A-\lambda I)M\bigr)=\det(A-\lambda I).
\]
Similarity is a change of basis, so the message is that eigenvalues belong to
the underlying transformation rather than to the particular array of numbers
used to record it --- a point taken up properly in
\refch{linear-transformations}.

\setchapterpreamble[u]{\margintoc}
\chapter{The Singular Value Decomposition}
\labch{svd}

\sources{Strang, \emph{Introduction to Linear Algebra}, 6th ed.\ \cite{strang2023ila},
§7.1--7.4; MIT 18.06 \cite{ocw1806}, Lecture~31.}

Eigenvalues require a square matrix and may fail to supply a full basis. The
singular value decomposition has neither restriction: every matrix has one. It
is the factorisation that ties the four subspaces together and the one that
underlies most of \refch{learning-from-data}.

\section{The statement}
\labsec{svd-statement}

\begin{theorem}
Every \(A\in\R^{m\times n}\) can be written
\[
	A=U\Sigma V\tp,
\]
where \(U\in\R^{m\times m}\) and \(V\in\R^{n\times n}\) are orthogonal and
\(\Sigma\in\R^{m\times n}\) is diagonal with entries
\(\sigma_1\ge\sigma_2\ge\cdots\ge\sigma_r>0\) and zeros elsewhere, \(r\) being
the rank of \(A\).
\end{theorem}

The columns \(\vect{v}_j\) are the \emph{right singular vectors}, the columns
\(\vect{u}_i\) the \emph{left singular vectors}, and the \(\sigma_i\) the
\emph{singular values}. The content of the factorisation is the pairing
\[
	A\vect{v}_i=\sigma_i\vect{u}_i,
	\qquad i=1,\ldots,r,
\]
which says that \(A\) maps one orthonormal set to another, up to individual
stretching factors.

\marginnote{Compare with eigenvectors, where the same basis appears on both
sides. Allowing two different bases is exactly what removes the need for \(A\)
to be square or diagonalisable.}

\section{Where it comes from}
\labsec{svd-construction}

Suppose \(A=U\Sigma V\tp\). Then
\[
	A\tp A=V\Sigma\tp U\tp U\Sigma V\tp=V(\Sigma\tp\Sigma)V\tp,
	\qquad
	AA\tp=U(\Sigma\Sigma\tp)U\tp .
\]
Both are symmetric and positive semidefinite, so \refsec{spectral} applies. The
\(\vect{v}_i\) are eigenvectors of \(A\tp A\), the \(\vect{u}_i\) are
eigenvectors of \(AA\tp\), and the singular values are the square roots of the
shared non-zero eigenvalues.

\begin{example}
Let \(A=\begin{bmatrix}3&0\\4&5\end{bmatrix}\). Then
\(A\tp A=\begin{bmatrix}25&20\\20&25\end{bmatrix}\), with eigenvalues \(45\) and
\(5\) and eigenvectors \([1,1]\tp\) and \([1,-1]\tp\). Hence
\(\sigma_1=3\sqrt5\), \(\sigma_2=\sqrt5\), and
\[
	\vect{v}_1=\tfrac{1}{\sqrt2}\begin{bmatrix}1\\1\end{bmatrix},
	\quad
	\vect{u}_1=\frac{A\vect{v}_1}{\sigma_1}=\tfrac{1}{\sqrt{10}}\begin{bmatrix}1\\3\end{bmatrix},
	\quad
	\vect{u}_2=\frac{A\vect{v}_2}{\sigma_2}=\tfrac{1}{\sqrt{10}}\begin{bmatrix}3\\-1\end{bmatrix}.
\]
As a check, \(\vect{u}_1\tp\vect{u}_2=0\) and
\(\sigma_1\sigma_2=3\sqrt5\cdot\sqrt5=15=\det A\).
\end{example}

\begin{remark}
In floating point one does not form \(A\tp A\); squaring the matrix squares its
condition number. The derivation above explains what the singular vectors are;
production algorithms compute them by orthogonal transformations applied to
\(A\) directly.
\end{remark}

\section{The four subspaces, one picture}
\labsec{svd-subspaces}

Split the singular vectors at index \(r\):
\[
	\underbrace{\vect{v}_1,\ldots,\vect{v}_r}_{\text{row space}},\quad
	\underbrace{\vect{v}_{r+1},\ldots,\vect{v}_n}_{\Null(A)},\qquad
	\underbrace{\vect{u}_1,\ldots,\vect{u}_r}_{\Col(A)},\quad
	\underbrace{\vect{u}_{r+1},\ldots,\vect{u}_m}_{\Null(A\tp)} .
\]
Each group is an orthonormal basis of the subspace beneath it. The orthogonality
statements of \refch{orthogonality} are then read off directly, and the action
of \(A\) becomes: discard the nullspace component, and stretch the surviving
\(r\) directions by \(\sigma_1,\ldots,\sigma_r\).

\section{Rank-one pieces and best approximation}
\labsec{eckart-young}

The column-times-row reading of \refsec{four-ways} turns the factorisation into
a sum ordered by importance:
\begin{equation}
	A=\sum_{i=1}^{r}\sigma_i\,\vect{u}_i\vect{v}_i\tp .
	\labeq{svd-sum}
\end{equation}

Truncating after \(k\) terms gives \(A_k\), and no other matrix of rank \(k\)
comes closer.

\begin{theorem}[Eckart--Young]
Among all matrices \(B\) of rank at most \(k\), the truncation \(A_k\) minimises
\(\lVert A-B\rVert\) in both the spectral and the Frobenius norm, with
\(\lVert A-A_k\rVert_2=\sigma_{k+1}\).
\end{theorem}

\marginnote{This single theorem is the justification for principal component
analysis, low-rank compression of weight matrices, and latent-factor models
alike. They differ in what the matrix contains, not in the mathematics applied
to it.}

\section{The pseudoinverse}
\labsec{pseudoinverse}

When \(A\) is not invertible, the closest available substitute is
\[
	A^{+}=V\Sigma^{+}U\tp,
\]
where \(\Sigma^{+}\) transposes \(\Sigma\) and replaces each \(\sigma_i>0\) by
\(1/\sigma_i\), leaving the zeros alone. Then \(\vect{x}=A^{+}\vect{b}\) is the
least-squares solution of smallest length: among all vectors minimising
\(\lVert A\vect{x}-\vect{b}\rVert\), it is the one lying entirely in the row
space.

\begin{remark}
This closes a loop opened in \refch{orthogonality}. There the normal equations
needed independent columns so that \(A\tp A\) could be inverted. The
pseudoinverse removes that condition and picks a canonical answer from the
infinitely many that rank deficiency allows.
\end{remark}

\setchapterpreamble[u]{\margintoc}
\chapter{Linear Transformations and Change of Basis}
\labch{linear-transformations}

\sources{Strang, \emph{Introduction to Linear Algebra}, 6th ed.\ \cite{strang2023ila},
§8.1--8.3; MIT 18.06 \cite{ocw1806}, Lectures~33--34.}

A matrix is a table of numbers; a linear transformation is a rule. This chapter
separates the two, which explains why the same rule can be recorded by many
different matrices and why choosing the basis well is worth the trouble.

\section{What makes a map linear}
\labsec{linear-maps}

\begin{definition}
A map \(T:V\to W\) between vector spaces is \emph{linear} when
\[
	T(\vect{u}+\vect{v})=T(\vect{u})+T(\vect{v}),
	\qquad
	T(c\vect{v})=cT(\vect{v})
\]
for all \(\vect{u},\vect{v}\in V\) and scalars \(c\).
\end{definition}

Setting \(c=0\) gives \(T(\vect{0})=\vect{0}\), so any map that shifts the origin
is disqualified. Rotation about the origin, projection onto a subspace, and
differentiation of polynomials are linear; translation and the map
\(\vect{v}\mapsto\lVert\vect{v}\rVert\) are not.

\marginnote{The two conditions combine into one: \(T\) is linear exactly when it
commutes with taking linear combinations. Everything in this book is a
consequence of that single permission.}

\section{The matrix of a transformation}
\labsec{transformation-matrix}

Fix a basis \(\vect{v}_1,\ldots,\vect{v}_n\) of \(V\) and a basis
\(\vect{w}_1,\ldots,\vect{w}_m\) of \(W\). Because \(T\) respects combinations,
it is determined by what it does to the basis vectors. Writing each image in the
output basis,
\[
	T(\vect{v}_j)=\sum_{i=1}^{m}a_{ij}\vect{w}_i,
\]
collects the coefficients into a matrix \(A\in\R^{m\times n}\), and the rule
becomes matrix multiplication in coordinates.

\begin{remark}
Column \(j\) of \(A\) is nothing other than the coordinate vector of
\(T(\vect{v}_j)\). Building the matrix of a transformation is always the same
mechanical act: feed in each basis vector, write the answer in the output basis,
and stack the results as columns.
\end{remark}

\begin{example}
Let \(V\) be the polynomials of degree at most \(3\), with basis
\(1,t,t^{2},t^{3}\), and let \(W\) be those of degree at most \(2\), with basis
\(1,t,t^{2}\). Differentiation sends \(1\mapsto0\), \(t\mapsto1\),
\(t^{2}\mapsto2t\), \(t^{3}\mapsto3t^{2}\), so
\[
	A=\begin{bmatrix}0&1&0&0\\0&0&2&0\\0&0&0&3\end{bmatrix}.
\]
Its nullspace is spanned by \([1,0,0,0]\tp\), the constants --- exactly the
polynomials that differentiate to zero --- and its rank is \(3\), matching
\(3+1=4\) in \refch{four-subspaces}.
\end{example}

\begin{example}
Rotation of \(\R^{2}\) by an angle \(\theta\) sends \(\vect{e}_1\) to
\((\cos\theta,\sin\theta)\) and \(\vect{e}_2\) to \((-\sin\theta,\cos\theta)\), so
\[
	A=\begin{bmatrix}\cos\theta&-\sin\theta\\ \sin\theta&\cos\theta\end{bmatrix},
\]
which satisfies \(A\tp A=I\): rotations preserve length, as an orthogonal matrix
must.
\end{example}

\section{Change of basis}
\labsec{change-of-basis}

Keeping the transformation and changing the basis changes the matrix. Let \(M\)
have as its columns the new basis vectors expressed in the old basis. A vector
with new coordinates \(\vect{c}\) has old coordinates \(M\vect{c}\), and the two
matrices of the same transformation are related by
\[
	A_{\text{new}}=M^{-1}A_{\text{old}}M .
\]

This is the similarity relation of \refsec{similar}, now with its meaning
supplied: similar matrices are the same transformation described by two
observers using different rulers. Eigenvalues, trace, determinant and rank are
properties of the transformation and are therefore identical for all of them;
the individual entries are not.

\marginnote{A useful sanity check when a formula looks wrong: ask whether the
quantity you computed could possibly depend on the choice of basis. If it
should not, and your expression does, the error is in the bookkeeping.}

\section{Choosing the basis well}
\labsec{good-basis}

Diagonalisation is the statement that, for a diagonalisable \(A\), the
eigenvector basis makes the matrix diagonal:
\(\Lambda=X^{-1}AX\). In that basis the transformation is \(n\) independent
scalings, and questions about powers, exponentials, and long-run behaviour become
questions about numbers.

When \(A\) is not square, or not diagonalisable, the singular value
decomposition supplies the substitute. Taking \(\vect{v}_1,\ldots,\vect{v}_n\) as
the input basis and \(\vect{u}_1,\ldots,\vect{u}_m\) as the output basis, the
matrix of the transformation is \(\Sigma\) --- diagonal, non-negative, and
ordered. No two-basis choice does better.

\begin{remark}
Read this way, the factorisations of Part~I are all answers to one question:
\emph{which coordinates make this transformation simple?} Elimination answers it
for solving, Gram--Schmidt for projecting, diagonalisation for iterating, and
the singular value decomposition for approximating.
\end{remark}

\setchapterpreamble[u]{\margintoc}
\chapter{Linear Algebra in Optimisation}
\labch{optimization}

\sources{Strang, \emph{Introduction to Linear Algebra}, 6th ed.\ \cite{strang2023ila},
Chapter~9.}

Optimisation problems that can be solved exactly are, with few exceptions,
quadratic ones. This chapter collects the linear algebra that governs them: the
norms used to measure error, the condition number that predicts difficulty, and
the quadratic form whose minimiser is a linear solve in disguise. The calculus
and the algorithms come in \refch{gradient-methods}.

\section{Norms}
\labsec{norms}

For vectors the three standard choices are
\[
	\lVert\vect{x}\rVert_1=\sum_i|x_i|,
	\qquad
	\lVert\vect{x}\rVert_2=\Bigl(\sum_i x_i^{2}\Bigr)^{1/2},
	\qquad
	\lVert\vect{x}\rVert_\infty=\max_i|x_i| .
\]
All three satisfy the triangle inequality; only \(\lVert\cdot\rVert_2\) comes
from an inner product, which is why projection and least squares are stated in
it.

\marginnote{The choice is not cosmetic. Penalising \(\lVert\cdot\rVert_1\)
drives coefficients exactly to zero and produces sparse solutions; penalising
\(\lVert\cdot\rVert_2\) shrinks them smoothly and keeps them non-zero.}

For matrices, two norms are used constantly:
\[
	\lVert A\rVert_2=\sigma_1,
	\qquad
	\lVert A\rVert_F=\Bigl(\sum_{i,j}a_{ij}^{2}\Bigr)^{1/2}
	=\Bigl(\sum_i\sigma_i^{2}\Bigr)^{1/2},
\]
the largest singular value and the Frobenius norm. Both are supplied by the
decomposition of \refch{svd}; the second identity says the Frobenius norm does
not care which orthonormal basis the matrix is written in.

\section{Conditioning}
\labsec{conditioning}

\begin{definition}
The \emph{condition number} of an invertible \(A\) is
\(\kappa(A)=\sigma_1/\sigma_n\), the ratio of largest to smallest singular value.
\end{definition}

It bounds how much a relative perturbation of \(\vect{b}\) can be amplified in
the solution of \(A\vect{x}=\vect{b}\): a relative change of size \(\varepsilon\)
in the data produces a relative change of at most \(\kappa(A)\varepsilon\) in the
answer. A matrix with \(\kappa=10^{8}\) loses roughly eight digits, whatever
algorithm is used.

\begin{remark}
This is the precise reason \refch{orthogonality} warned against forming
\(A\tp A\). Its singular values are the squares of those of \(A\), so
\(\kappa(A\tp A)=\kappa(A)^{2}\) and the difficulty is squared along with them.
\end{remark}

\section{Quadratic forms}
\labsec{quadratic-forms}

A quadratic function on \(\R^{n}\) can always be written
\[
	f(\vect{x})=\tfrac12\vect{x}\tp A\vect{x}-\vect{b}\tp\vect{x}+c,
\]
with \(A\) symmetric --- any non-symmetric part cancels in the product and can be
discarded without changing \(f\).

\begin{theorem}
If \(A\) is symmetric positive definite then \(f\) has a unique minimiser, and it
is the solution of \(A\vect{x}=\vect{b}\).
\end{theorem}

\begin{proof}
Let \(A\vect{x}^{\star}=\vect{b}\) and write \(\vect{x}=\vect{x}^{\star}+\vect{d}\).
Expanding and using symmetry,
\[
	f(\vect{x})=f(\vect{x}^{\star})+\tfrac12\vect{d}\tp A\vect{d},
\]
and positive definiteness makes the last term strictly positive unless
\(\vect{d}=\vect{0}\).
\end{proof}

\marginnote{Minimisation and solving are therefore the same problem wearing two
hats. Least squares in \refch{orthogonality} was the first instance:
\(\lVert A\vect{x}-\vect{b}\rVert^{2}\) is a quadratic form whose matrix is
\(A\tp A\).}

The level sets of \(f\) are ellipsoids whose axes are the eigenvectors of \(A\),
with lengths proportional to \(1/\sqrt{\lambda_i}\). A large ratio
\(\lambda_{\max}/\lambda_{\min}\) makes them long and thin, which is exactly the
geometry that slows gradient methods down in \refch{gradient-methods}.

\section{Regularisation}
\labsec{regularisation}

When \(A\tp A\) is singular or nearly so, adding a multiple of the identity
repairs it:
\[
	(A\tp A+\delta^{2}I)\hat{\vect{x}}=A\tp\vect{b},
	\qquad \delta>0 .
\]
The matrix on the left is positive definite for every \(\delta>0\), since
\(\vect{x}\tp(A\tp A+\delta^{2}I)\vect{x}
=\lVert A\vect{x}\rVert^{2}+\delta^{2}\lVert\vect{x}\rVert^{2}\). In terms of the
singular value decomposition each \(\sigma_i\) is replaced by
\(\sigma_i/(\sigma_i^{2}+\delta^{2})\), so small singular values are damped
rather than inverted.

\begin{remark}
Statistics calls this ridge regression and deep learning calls it weight decay,
but the linear algebra is identical: trade a little bias for a large reduction
in sensitivity to noise. As \(\delta\to0\) the solution approaches the
pseudoinverse answer of \refsec{pseudoinverse}.
\end{remark}

\setchapterpreamble[u]{\margintoc}
\chapter{Linear Algebra behind Learning from Data}
\labch{learning-from-data}

\sources{Strang, \emph{Introduction to Linear Algebra}, 6th ed.\ \cite{strang2023ila},
Chapter~10; Zhang, Lipton, Li and Smola, \emph{Dive into Deep Learning}
\cite{d2l2023}, Chapter~2.}

This chapter closes Part~I by putting the factorisations to work on data. Nothing
new is proved; the point is to see that the objects of the previous nine
chapters are the objects a learning system manipulates.

\section{Data as a matrix}
\labsec{data-matrix}

Store \(n\) examples with \(d\) features each as \(X\in\R^{n\times d}\), one
example per row. Two readings, both from \refsec{four-ways}, are used
constantly:

\begin{description}
	\item[By rows] each row is a point in \(\R^{d}\), and the dataset is a cloud
	of \(n\) such points.
	\item[By columns] each column is one feature measured across all examples,
	and the column space of \(X\) is the set of quantities expressible as
	weighted mixtures of features.
\end{description}

A linear model is then \(\hat{\vect{y}}=X\vect{w}\), so the achievable
predictions form exactly \(\Col(X)\), and fitting it by least squares is the
projection of \refch{orthogonality}.

\marginnote{This is the sense in which linear algebra is not a prerequisite for
learning but a description of it. The reachability question of
\refch{vectors-and-column-space} is the question of what a linear model can
represent.}

\section{Centring and covariance}
\labsec{covariance}

Subtracting the mean of each column gives the centred matrix \(X_0\). The
\emph{sample covariance matrix} is
\[
	S=\frac{1}{n-1}X_0\tp X_0\in\R^{d\times d},
\]
symmetric and positive semidefinite by \refsec{positive-definite}, with
\(s_{jj}\) the variance of feature \(j\) and \(s_{jk}\) the covariance between
features \(j\) and \(k\).

\begin{example}
Let
\[
	X_0=\begin{bmatrix}2&1\\1&2\\-2&-1\\-1&-2\end{bmatrix},
\]
whose columns already sum to zero. Then
\(X_0\tp X_0=\begin{bmatrix}10&8\\8&10\end{bmatrix}\) and
\(S=\tfrac13X_0\tp X_0\). The eigenvalues of \(X_0\tp X_0\) are \(18\) and \(2\),
with eigenvectors \([1,1]\tp\) and \([1,-1]\tp\).
\end{example}

\section{Principal components}
\labsec{pca}

Principal component analysis asks for the direction \(\vect{q}\) of unit length
along which the projected data vary most. That is the maximiser of
\(\vect{q}\tp S\vect{q}\), and by the spectral theorem it is the eigenvector of
\(S\) belonging to the largest eigenvalue.

\begin{proposition}
The principal directions of \(X_0\) are its right singular vectors, and the
variance captured by the \(i\)th component is proportional to \(\sigma_i^{2}\).
\end{proposition}

\begin{proof}
Writing \(X_0=U\Sigma V\tp\) gives \(X_0\tp X_0=V\Sigma\tp\Sigma V\tp\), which is
the eigendecomposition of a symmetric matrix with eigenvalues \(\sigma_i^{2}\)
and eigenvectors \(\vect{v}_i\).
\end{proof}

\begin{example}[continued]
Here \(\sigma_1=\sqrt{18}=3\sqrt2\) and \(\sigma_2=\sqrt2\), so the first
component \([1,1]\tp/\sqrt2\) accounts for \(18/20=90\%\) of the total variance.
Projecting onto it replaces two features by one while discarding a tenth of the
spread.
\end{example}

\marginnote{Computing the singular value decomposition of \(X_0\) is preferable
to forming \(S\) first, for the conditioning reason given in
\refsec{conditioning}.}

\section{Low rank as a modelling assumption}
\labsec{low-rank}

Eckart--Young (\refsec{eckart-young}) says the truncated sum
\(A_k=\sum_{i\le k}\sigma_i\vect{u}_i\vect{v}_i\tp\) is the best rank-\(k\)
approximation. Three familiar techniques are that theorem applied to different
matrices: dimensionality reduction applies it to a data matrix, compression of a
trained layer applies it to a weight matrix, and latent-factor recommendation
applies it to a partially observed ratings matrix.

\begin{remark}
The assumption being made each time is that the interesting structure lives in
few directions and the rest is noise. Whether that is true is a modelling
question about the data, not a mathematical one --- the theorem is silent about
where to cut.
\end{remark}

\section{Batched computation}
\labsec{batching}

A linear layer applied to a batch is one matrix product. With \(X\in\R^{n\times d}\)
and weights \(W\in\R^{d\times h}\),
\[
	Z=XW+\vect{1}\vect{b}\tp\in\R^{n\times h},
\]
where \(\vect{1}\in\R^{n}\) broadcasts the bias across rows. The block reading of
\refsec{four-ways} is what makes this efficient: the product decomposes into
independent blocks that hardware evaluates in parallel, which is the practical
reason mini-batches exist at all.

\begin{remark}
Bias makes the layer \emph{affine} rather than linear, so it violates
\(T(\vect{0})=\vect{0}\) from \refsec{linear-maps}. Nothing breaks --- an affine
map is a linear map followed by a shift --- but statements proved for linear
maps have to be applied to the linear part alone.
\end{remark}

\section{What carries forward}
\labsec{part1-summary}

Four ideas from Part~I are used without further comment in the rest of the book:
the column space as the set of reachable outputs; orthogonal projection as the
solution of a least-squares problem; the eigenvalues of a symmetric matrix as
the curvature of a quadratic; and the singular value decomposition as the
canonical way to separate signal from remainder. \refch{column-space-to-layers}
returns to them once the calculus and probability are in place.


\pagelayout{wide}
\addpart{Calculus for Learning}
\pagelayout{margin}

\setchapterpreamble[u]{\margintoc}
\chapter{Derivatives and the Chain Rule}
\labch{derivatives}

\sources{Hass, Heil and Weir, \emph{Thomas' Calculus}, 14th ed.\ \cite{thomas2018calculus},
Chapters~2--3; MIT 18.01 \cite{ocw1801}.}

Part~II covers only the calculus that optimisation needs. That means
differentiation, and in particular the chain rule, which is the entire
mathematical content of training a network by gradients. Integration appears
only where \refch{expectation} needs it.

\section{The derivative as a local linear approximation}
\labsec{derivative-definition}

\begin{definition}
For \(f:\R\to\R\), the derivative at \(a\) is
\[
	f'(a)=\lim_{h\to0}\frac{f(a+h)-f(a)}{h},
\]
when the limit exists.
\end{definition}

The definition is a statement about approximation. Rearranged, it says
\[
	f(a+h)=f(a)+f'(a)h+o(h),
\]
where the remainder shrinks faster than \(h\). In words: near \(a\), the function
is the straight line \(f(a)+f'(a)h\), plus an error too small to matter at first
order.

\marginnote{This reading is the one that generalises. In \refch{gradients} the
straight line becomes a plane and \(f'(a)\) becomes a vector; in
\refch{jacobian-hessian} the plane becomes an affine map and the vector becomes
a matrix. The idea never changes.}

Differentiability implies continuity, but not conversely: \(f(x)=|x|\) is
continuous at \(0\) and has no derivative there, because the one-sided limits
disagree. This matters in practice, since the rectified linear unit
\(\max(0,x)\) has exactly this kink and is handled by convention rather than by
theorem.

\section{The rules}
\labsec{diff-rules}

For differentiable \(f\) and \(g\),
\[
	(f+g)'=f'+g',
	\qquad
	(fg)'=f'g+fg',
	\qquad
	\Bigl(\frac{f}{g}\Bigr)'=\frac{f'g-fg'}{g^{2}},
\]
together with \((x^{n})'=nx^{n-1}\), \((e^{x})'=e^{x}\), and
\((\ln x)'=1/x\) for \(x>0\).

\begin{example}
The logistic function \(\sigma(z)=\dfrac{1}{1+e^{-z}}\) satisfies
\[
	\sigma'(z)=\frac{e^{-z}}{(1+e^{-z})^{2}}=\sigma(z)\bigl(1-\sigma(z)\bigr).
\]
The derivative is expressible in terms of the value itself, which is why
implementations cache \(\sigma(z)\) during the forward pass and reuse it going
backwards.
\end{example}

\section{The chain rule}
\labsec{chain-rule}

\begin{theorem}
If \(g\) is differentiable at \(a\) and \(f\) is differentiable at \(g(a)\), then
\(f\circ g\) is differentiable at \(a\) and
\[
	(f\circ g)'(a)=f'\bigl(g(a)\bigr)\,g'(a).
\]
\end{theorem}

The linear-approximation reading makes the statement almost obvious: \(g\)
multiplies a small input change by \(g'(a)\), then \(f\) multiplies the result by
\(f'(g(a))\), so the composed effect is the product of the two factors.

\begin{example}
For \(f(x)=\sigma(wx+b)\) with \(w,b\) fixed,
\(f'(x)=\sigma'(wx+b)\cdot w\). Differentiating instead with respect to the
\emph{parameter} \(w\) gives \(\sigma'(wx+b)\cdot x\). Same rule, different
variable held fixed --- and it is the second version that training uses.
\end{example}

\marginnote{Almost every gradient in a neural network is this calculation
repeated. What makes networks distinctive is not the rule but the number of
times it is applied and the care taken over the order of the arithmetic.}

\section{Long compositions}
\labsec{long-compositions}

A network of \(L\) layers is a composition
\(f=f_L\circ f_{L-1}\circ\cdots\circ f_1\), and the chain rule gives
\[
	f'(x)=f_L'\bigl(z_{L-1}\bigr)\cdot f_{L-1}'\bigl(z_{L-2}\bigr)\cdots f_1'(x),
\]
where \(z_k\) is the output of layer \(k\). Two consequences deserve naming even
in the scalar case.

First, the derivative is a \emph{product} of \(L\) factors. If the factors are
systematically smaller than one the product decays geometrically; if
systematically larger it explodes. This is the vanishing and exploding gradient
problem, visible here with no matrices involved.

Second, the intermediate values \(z_k\) are needed to evaluate the factors. They
are computed on the way forward and must be kept, which is why memory during
training scales with depth.

\begin{remark}
Multiplication is associative but not equally cheap in every order. Accumulating
the product from the output end backwards keeps every intermediate quantity the
same shape as the output; accumulating from the input end does not.
\refsec{backprop} makes this precise once the objects are matrices.
\end{remark}

\section{What is not covered}
\labsec{calculus-scope}

Techniques of integration, infinite series, and differential equations are
outside the scope agreed for these notes. Definite integrals appear in
\refch{expectation} only to the extent that expectations of continuous random
variables require them, and the fundamental theorem of calculus is used there
without a separate development.

\setchapterpreamble[u]{\margintoc}
\chapter{Partial Derivatives and Gradients}
\labch{gradients}

\sources{Hass, Heil and Weir, \emph{Thomas' Calculus}, 14th ed.\ \cite{thomas2018calculus},
Chapter~14; MIT 18.02 \cite{ocw1801}.}

A loss depends on many parameters at once, so the derivative has to become a
vector. This chapter builds that vector and computes it for the handful of
expressions that account for most of what appears in practice.

\section{Partial derivatives}
\labsec{partials}

For \(f:\R^{n}\to\R\), the partial derivative \(\partial f/\partial x_i\) is the
ordinary derivative with respect to \(x_i\) with the other coordinates held
fixed. Nothing new is required to compute one; the novelty is in assembling
them.

\begin{definition}
The \emph{gradient} of \(f\) at \(\vect{x}\) is the column vector
\[
	\nabla f(\vect{x})
	=\begin{bmatrix}
		\partial f/\partial x_1\\ \vdots\\ \partial f/\partial x_n
	\end{bmatrix}.
\]
\end{definition}

The first-order approximation of \refsec{derivative-definition} becomes
\begin{equation}
	f(\vect{x}+\vect{h})=f(\vect{x})+\nabla f(\vect{x})\tp\vect{h}+o(\lVert\vect{h}\rVert),
	\labeq{first-order}
\end{equation}
so the gradient is the vector whose inner product with a displacement predicts
the change in \(f\).

\section{Direction of steepest ascent}
\labsec{steepest}

\begin{proposition}
Among unit vectors \(\vect{u}\), the directional derivative
\[
	\nabla f(\vect{x})\tp\vect{u}
\]
is largest when \(\vect{u}\) points along \(\nabla f(\vect{x})\), and its value
there is \(\lVert\nabla f(\vect{x})\rVert\).
\end{proposition}

\begin{proof}
By the Cauchy--Schwarz inequality,
\(|\nabla f\tp\vect{u}|\le\lVert\nabla f\rVert\lVert\vect{u}\rVert
=\lVert\nabla f\rVert\), with equality exactly when \(\vect{u}\) is parallel to
\(\nabla f\).
\end{proof}

\marginnote{Hence \(-\nabla f\) is the direction of steepest \emph{descent}. That
single sentence is the justification for every method in
\refch{gradient-methods}.}

A related fact is that the gradient is perpendicular to the level set through
\(\vect{x}\): moving along a level set does not change \(f\), so the inner
product in \refeq{first-order} must vanish for such displacements.

\section{Gradients of the expressions that matter}
\labsec{matrix-gradients}

Three formulas cover most cases encountered in this book. Throughout,
\(\vect{a}\) is a fixed vector and \(A\) a fixed matrix.

\begin{proposition}
\begin{align*}
	\nabla\bigl(\vect{a}\tp\vect{x}\bigr)&=\vect{a},\\
	\nabla\bigl(\vect{x}\tp A\vect{x}\bigr)&=(A+A\tp)\vect{x},\\
	\nabla\Bigl(\tfrac12\lVert A\vect{x}-\vect{b}\rVert^{2}\Bigr)
		&=A\tp(A\vect{x}-\vect{b}).
\end{align*}
In particular the middle formula is \(2A\vect{x}\) when \(A\) is symmetric.
\end{proposition}

\begin{proof}[Third formula]
Write \(f(\vect{x})=\tfrac12(A\vect{x}-\vect{b})\tp(A\vect{x}-\vect{b})\) and
expand at \(\vect{x}+\vect{h}\):
\[
	f(\vect{x}+\vect{h})=f(\vect{x})
	+(A\vect{x}-\vect{b})\tp A\vect{h}
	+\tfrac12\lVert A\vect{h}\rVert^{2}.
\]
The middle term is \(\bigl(A\tp(A\vect{x}-\vect{b})\bigr)\tp\vect{h}\) and the
last is \(O(\lVert\vect{h}\rVert^{2})\); comparing with \refeq{first-order}
identifies the gradient.
\end{proof}

\begin{remark}
Setting the third gradient to zero gives \(A\tp A\vect{x}=A\tp\vect{b}\) --- the
normal equations of \refeq{normal-equations}. The geometric derivation in
\refch{orthogonality} and the calculus derivation here produce the same
equations, as they must.
\end{remark}

\section{A worked parameter gradient}
\labsec{worked-gradient}

Take one logistic unit with squared loss on a single example \((\vect{x},y)\):
\[
	z=\vect{w}\tp\vect{x}+b,
	\qquad
	\hat y=\sigma(z),
	\qquad
	\ell=\tfrac12(\hat y-y)^{2}.
\]
Working outwards from the loss,
\[
	\frac{\partial\ell}{\partial\hat y}=\hat y-y,
	\qquad
	\frac{\partial\hat y}{\partial z}=\hat y(1-\hat y),
	\qquad
	\frac{\partial z}{\partial\vect{w}}=\vect{x},
\]
so that
\[
	\nabla_{\vect{w}}\ell=(\hat y-y)\,\hat y(1-\hat y)\,\vect{x},
	\qquad
	\frac{\partial\ell}{\partial b}=(\hat y-y)\,\hat y(1-\hat y).
\]

\marginnote{Note the factor \(\hat y(1-\hat y)\), which is near zero whenever the
unit is confidently right or confidently wrong. Squared loss on a saturating
unit therefore learns slowly from its worst mistakes --- one of the reasons
cross-entropy is preferred, as \refch{estimation} explains.}

\section{Convexity in one line}
\labsec{convexity}

A differentiable \(f\) is \emph{convex} when
\[
	f(\vect{y})\ge f(\vect{x})+\nabla f(\vect{x})\tp(\vect{y}-\vect{x})
	\qquad\text{for all }\vect{x},\vect{y},
\]
that is, when the first-order approximation never overestimates. For convex
\(f\) a point with \(\nabla f=\vect{0}\) is a global minimum, which is why
convex problems are the ones that can be declared solved. The losses of interest
in Part~IV are not convex in the parameters, and \refch{gradient-methods} says
what survives of the theory anyway.

\setchapterpreamble[u]{\margintoc}
\chapter{Jacobians, Hessians, and Taylor Approximation}
\labch{jacobian-hessian}

\sources{Hass, Heil and Weir, \emph{Thomas' Calculus}, 14th ed.\ \cite{thomas2018calculus},
Chapters~14--15; MIT 18.02 \cite{ocw1801}.}

The gradient handles scalar-valued functions. Vector-valued functions need a
matrix of derivatives, and second-order behaviour needs another. Both are
already familiar objects from Part~I wearing new labels.

\section{The Jacobian}
\labsec{jacobian}

\begin{definition}
For \(\vect{f}:\R^{n}\to\R^{m}\), the \emph{Jacobian} at \(\vect{x}\) is
\(J\in\R^{m\times n}\) with entries
\[
	J_{ij}=\frac{\partial f_i}{\partial x_j}.
\]
\end{definition}

Row \(i\) of \(J\) is the gradient of the \(i\)th output. The first-order
approximation is
\[
	\vect{f}(\vect{x}+\vect{h})=\vect{f}(\vect{x})+J\vect{h}+o(\lVert\vect{h}\rVert),
\]
so the Jacobian is the matrix of the linear map that best imitates \(\vect{f}\)
near \(\vect{x}\). When \(m=1\) it is \(\nabla f\tp\), a single row.

\marginnote{Note the transposition. The gradient is a column by convention, the
Jacobian of a scalar function a row. Keeping this straight is the source of most
sign-free errors in gradient derivations.}

\begin{example}
An affine layer \(\vect{f}(\vect{x})=W\vect{x}+\vect{b}\) has \(J=W\)
everywhere: it is already linear, so its best linear approximation is itself. An
elementwise activation \(\vect{f}(\vect{x})=\bigl(\phi(x_1),\ldots,\phi(x_n)\bigr)\)
has \(J=\operatorname{diag}\bigl(\phi'(x_1),\ldots,\phi'(x_n)\bigr)\).
\end{example}

\section{The chain rule in matrix form}
\labsec{matrix-chain-rule}

\begin{theorem}
If \(\vect{h}=\vect{g}\circ\vect{f}\) with \(\vect{f}:\R^{n}\to\R^{p}\) and
\(\vect{g}:\R^{p}\to\R^{m}\), then
\[
	J_{\vect{h}}(\vect{x})=J_{\vect{g}}\bigl(\vect{f}(\vect{x})\bigr)\,J_{\vect{f}}(\vect{x}).
\]
\end{theorem}

The scalar chain rule of \refsec{chain-rule} is the case \(m=p=n=1\). The
matrix version says a composition of maps corresponds to a product of Jacobians,
in the same order as the composition --- and since matrix multiplication is not
commutative, the order is now load-bearing.

\begin{remark}
A network is exactly such a composition, alternating affine maps and elementwise
activations. Its Jacobian is therefore an alternating product
\(W_L D_{L-1} W_{L-1}\cdots D_1 W_1\) with \(D_k\) diagonal. The vanishing and
exploding behaviour anticipated in \refsec{long-compositions} is now a statement
about the singular values of that product.
\end{remark}

\section{The Hessian}
\labsec{hessian}

\begin{definition}
For \(f:\R^{n}\to\R\) with continuous second partials, the \emph{Hessian} is
\(H\in\R^{n\times n}\) with \(H_{ij}=\partial^{2}f/\partial x_i\partial x_j\).
\end{definition}

Equality of mixed partials makes \(H\) symmetric, so the spectral theorem of
\refsec{spectral} applies: it has real eigenvalues and an orthonormal
eigenbasis. Those eigenvalues are the curvatures of \(f\) along the
corresponding directions.

\section{Taylor approximation}
\labsec{taylor}

Extending \refeq{first-order} one order further,
\begin{equation}
	f(\vect{x}+\vect{h})=f(\vect{x})+\nabla f(\vect{x})\tp\vect{h}
	+\tfrac12\,\vect{h}\tp H(\vect{x})\vect{h}+o(\lVert\vect{h}\rVert^{2}).
	\labeq{taylor2}
\end{equation}

The quadratic term is exactly the quadratic form of \refsec{quadratic-forms}, so
everything proved there transfers. In particular a positive definite Hessian
makes the local model a bowl with a unique minimum.

\begin{proposition}[Second-derivative test]
Let \(\nabla f(\vect{x})=\vect{0}\). If \(H(\vect{x})\) is positive definite the
point is a local minimum; if negative definite, a local maximum; if it has
eigenvalues of both signs, a saddle point.
\end{proposition}

\begin{example}
For \(f(x_1,x_2)=x_1^{2}+3x_1x_2+2x_2^{2}\),
\(H=\begin{bmatrix}2&3\\3&4\end{bmatrix}\) with \(\det H=-1<0\), so the
eigenvalues have opposite signs and the origin is a saddle.
\end{example}

\marginnote{In high dimensions saddles vastly outnumber local minima, since a
critical point is a minimum only if \emph{every} one of \(n\) eigenvalues is
positive. Much of the practical behaviour of gradient descent on deep models is
about escaping saddles rather than escaping minima.}

\section{The Jacobian determinant}
\labsec{jacobian-determinant}

When \(m=n\), \(|\det J|\) measures how the map scales volume locally --- the
statement anticipated at the end of \refch{determinants}. It appears whenever
variables are changed inside an integral:
\[
	\int_{\vect{g}(\Omega)}f(\vect{y})\,d\vect{y}
	=\int_{\Omega}f\bigl(\vect{g}(\vect{x})\bigr)\,\bigl|\det J_{\vect{g}}(\vect{x})\bigr|\,d\vect{x}.
\]

\begin{remark}
This is the one place Part~II needs integration, and it is needed for
\refch{random-variables}: transforming a random variable transforms its density
by exactly this factor. A map with \(\det J=0\) collapses volume, which is the
calculus reading of singularity.
\end{remark}

\setchapterpreamble[u]{\margintoc}
\chapter{Gradient-Based Optimisation}
\labch{gradient-methods}

\sources{Hass, Heil and Weir, \emph{Thomas' Calculus}, 14th ed.\ \cite{thomas2018calculus},
Chapter~14; Strang, \emph{Introduction to Linear Algebra}, 6th ed.\
\cite{strang2023ila}, Chapter~9; Zhang, Lipton, Li and Smola,
\emph{Dive into Deep Learning} \cite{d2l2023}, Chapters~11--12.}

This chapter closes Part~II by combining the two halves already assembled: the
gradient says which way to move, and the Hessian says how far the move can be
trusted.

\section{Descent directions and the basic method}
\labsec{descent}

By \refsec{steepest}, \(-\nabla f\) is the direction along which \(f\) falls
fastest, so the simplest possible method is
\begin{equation}
	\vect{x}_{k+1}=\vect{x}_k-\alpha\,\nabla f(\vect{x}_k),
	\labeq{gd}
\end{equation}
with step size \(\alpha>0\). Everything interesting is in the choice of
\(\alpha\), and \refeq{taylor2} explains why: the linear model that justified the
direction is only accurate while the quadratic term stays small.

\section{The quadratic case, exactly}
\labsec{gd-quadratic}

Take \(f(\vect{x})=\tfrac12\vect{x}\tp A\vect{x}-\vect{b}\tp\vect{x}\) with \(A\)
symmetric positive definite, whose minimiser \(\vect{x}^{\star}\) solves
\(A\vect{x}=\vect{b}\) by \refsec{quadratic-forms}. Since
\(\nabla f(\vect{x})=A\vect{x}-\vect{b}=A(\vect{x}-\vect{x}^{\star})\), the error
\(\vect{e}_k=\vect{x}_k-\vect{x}^{\star}\) obeys
\[
	\vect{e}_{k+1}=(I-\alpha A)\,\vect{e}_k .
\]

Diagonalising \(A=Q\Lambda Q\tp\) decouples the coordinates: along the
eigenvector \(\vect{q}_i\) the error is multiplied by \(1-\alpha\lambda_i\) at
every step.

\begin{theorem}
Gradient descent on this quadratic converges for every starting point if and
only if \(0<\alpha<2/\lambda_{\max}\). The best rate is obtained at
\(\alpha^{\star}=2/(\lambda_{\min}+\lambda_{\max})\), where the error contracts
by a factor
\[
	\frac{\kappa-1}{\kappa+1},
	\qquad
	\kappa=\frac{\lambda_{\max}}{\lambda_{\min}},
\]
per step.
\end{theorem}

\marginnote{With \(\kappa=100\) the factor is \(0.980\); reducing the error by
\(10^{-3}\) then takes about \(340\) steps. Conditioning, introduced in
\refsec{conditioning} as a numerical concern, is here a direct statement about
training time.}

\begin{remark}
The geometry is the elongated ellipse of \refsec{quadratic-forms}. The negative
gradient points across the valley rather than along it, so the iterates zigzag.
Nothing about this depends on the objective being artificial --- an
ill-conditioned Hessian produces the same behaviour for any smooth loss.
\end{remark}

\section{Momentum and adaptive steps}
\labsec{momentum}

Momentum averages successive gradients,
\[
	\vect{v}_{k+1}=\beta\vect{v}_k+\nabla f(\vect{x}_k),
	\qquad
	\vect{x}_{k+1}=\vect{x}_k-\alpha\vect{v}_{k+1},
\]
which cancels the alternating components of the zigzag while accumulating the
consistent ones. Adaptive methods instead rescale each coordinate by a running
measure of its own gradient magnitude, which amounts to approximating the
diagonal of the Hessian on the fly.

\begin{remark}
Both are attempts to reduce the effective condition number without forming
second derivatives. Neither changes the fact, established above, that the
condition number is what has to be reduced.
\end{remark}

\section{Newton's method}
\labsec{newton}

Minimising the quadratic model \refeq{taylor2} exactly, rather than following its
linear part, gives
\[
	\vect{x}_{k+1}=\vect{x}_k-H(\vect{x}_k)^{-1}\nabla f(\vect{x}_k).
\]
On a quadratic objective this reaches the minimum in a single step, and near a
non-degenerate minimum it converges quadratically. The obstacle is cost: forming
and factorising \(H\) is \(O(n^{3})\) work in \(n\) parameters, which is
unavailable when \(n\) is in the millions. Quasi-Newton methods maintain a
cheap approximation instead.

\section{Stochastic gradients}
\labsec{sgd}

Training losses are sums over examples,
\(f(\vect{x})=\frac1n\sum_{i=1}^{n}\ell_i(\vect{x})\). Evaluating the full
gradient costs one pass over the data; evaluating it on a random mini-batch
costs a fraction of that and gives an unbiased estimate, since the expectation of
the batch gradient is the full gradient.

\begin{remark}
The trade is variance for speed. Because the estimate is noisy, a constant step
size leaves the iterates bouncing near the optimum rather than settling, which is
why step sizes are decayed. The unbiasedness claim itself belongs to
\refch{expectation}.
\end{remark}

\section{Backpropagation}
\labsec{backprop}

For a composition \(f=f_L\circ\cdots\circ f_1\) with a scalar output, the chain
rule of \refsec{matrix-chain-rule} gives
\[
	\nabla f=J_1\tp J_2\tp\cdots J_L\tp\,,
\]
a product that may be evaluated in either order. Multiplying from the left
(inputs first) propagates full Jacobians, each of size (layer width) squared.
Multiplying from the right (output first) propagates a single vector at every
stage, because the output is scalar.

Backpropagation is that second ordering, nothing more: a forward pass storing
the intermediate values, then a backward pass carrying one vector through the
transposed Jacobians. Its cost is a small constant multiple of the forward pass,
independent of the number of parameters, and that asymptotic fact is what makes
training large models possible at all.

\marginnote{The name suggests a special algorithm. It is better understood as a
choice of association order in a matrix product --- the choice that keeps every
intermediate object as small as the output rather than as large as the layer.}


\pagelayout{wide}
\addpart{Probability for Learning}
\pagelayout{margin}

\setchapterpreamble[u]{\margintoc}
\chapter{Sample Spaces, Conditioning, and Bayes}
\labch{sample-spaces}

\sources{Bertsekas and Tsitsiklis, \emph{Introduction to Probability}, 2nd ed.\
\cite{bertsekas2008probability}, Chapter~1; MIT 6.041 \cite{ocw6041},
Lectures~1--4; MIT 18.05 \cite{ocw1805}.}

Part~III covers the probability that a learning problem needs: enough to state
what a model assumes about its data, and enough to derive the two loss functions
that dominate practice. It stops at estimation, as agreed.

\section{Sample spaces and the axioms}
\labsec{axioms}

An experiment is described by a \emph{sample space} \(\Omega\), the set of
possible outcomes. An \emph{event} is a subset of \(\Omega\). A probability
measure assigns to each event a number \(\Pr(A)\) subject to three requirements:

\begin{enumerate}
	\item \(\Pr(A)\ge0\);
	\item \(\Pr(\Omega)=1\);
	\item \(\Pr\bigl(\bigcup_i A_i\bigr)=\sum_i\Pr(A_i)\) for disjoint events.
\end{enumerate}

\marginnote{Everything else is bookkeeping on top of these three lines ---
including the results in \refch{limit-theorems} that justify learning from
samples at all.}

Immediate consequences are \(\Pr(A^{c})=1-\Pr(A)\), monotonicity, and the
inclusion--exclusion formula
\(\Pr(A\cup B)=\Pr(A)+\Pr(B)-\Pr(A\cap B)\).

\section{Conditional probability}
\labsec{conditional-probability}

\begin{definition}
For \(\Pr(B)>0\), the probability of \(A\) given \(B\) is
\[
	\Pr(A\mid B)=\frac{\Pr(A\cap B)}{\Pr(B)}.
\]
\end{definition}

Conditioning is renormalisation: \(B\) becomes the new sample space, and the
measure is rescaled so that it again totals one. For fixed \(B\) the map
\(A\mapsto\Pr(A\mid B)\) satisfies the three axioms, so every result proved for
probabilities holds for conditional probabilities as well.

Rearranging gives the multiplication rule
\(\Pr(A\cap B)=\Pr(A\mid B)\Pr(B)\), and chaining it,
\[
	\Pr(A_1\cap\cdots\cap A_n)
	=\Pr(A_1)\Pr(A_2\mid A_1)\cdots\Pr(A_n\mid A_1\cap\cdots\cap A_{n-1}).
\]

\section{Total probability and Bayes' rule}
\labsec{bayes}

Let \(B_1,\ldots,B_k\) partition \(\Omega\). Splitting any event along the
partition gives the \emph{total probability} formula
\[
	\Pr(A)=\sum_{i=1}^{k}\Pr(A\mid B_i)\Pr(B_i),
\]
and dividing the multiplication rule by it gives \emph{Bayes' rule},
\begin{equation}
	\Pr(B_j\mid A)=\frac{\Pr(A\mid B_j)\Pr(B_j)}{\sum_{i}\Pr(A\mid B_i)\Pr(B_i)} .
	\labeq{bayes}
\end{equation}

\begin{example}
A test for a condition present in \(0.1\%\) of a population detects it in
\(99\%\) of carriers and returns a false positive for \(1\%\) of non-carriers.
For a positive result,
\[
	\Pr(\text{carrier}\mid+)
	=\frac{0.99\times0.001}{0.99\times0.001+0.01\times0.999}
	=\frac{0.00099}{0.01098}\approx0.090 .
\]
A test that is right \(99\%\) of the time leaves the subject about \(9\%\) likely
to be a carrier, because the prior is small enough that false positives
outnumber true ones by roughly ten to one.
\end{example}

\marginnote{The same arithmetic explains why a classifier with excellent accuracy
on a rare class can still produce mostly wrong positive predictions. Accuracy and
precision answer different questions; \refeq{bayes} is the conversion between
them.}

\section{Independence}
\labsec{independence}

\begin{definition}
Events \(A\) and \(B\) are \emph{independent} when
\(\Pr(A\cap B)=\Pr(A)\Pr(B)\), equivalently \(\Pr(A\mid B)=\Pr(A)\) when
\(\Pr(B)>0\).
\end{definition}

Independence is an assumption about the mechanism generating the data, not a
property that can be read off a formula. It is also fragile: events may be
pairwise independent without being jointly independent, and independence
generally fails once one conditions on a third event.

\begin{remark}
The i.i.d.\ assumption of \refch{learning-problem} is exactly independence
applied to the examples in a dataset, plus the requirement that they share one
distribution. Nearly every guarantee in \refch{limit-theorems} rests on it, and
distribution shift is the name for what happens when the second half fails.
\end{remark}

\setchapterpreamble[u]{\margintoc}
\chapter{Random Variables and Distributions}
\labch{random-variables}

\sources{Bertsekas and Tsitsiklis, \emph{Introduction to Probability}, 2nd ed.\
\cite{bertsekas2008probability}, Chapters~2--3; MIT 6.041 \cite{ocw6041},
Lectures~5--12.}

A random variable attaches a number to each outcome, which turns questions about
events into questions about functions --- and therefore into questions that
calculus and linear algebra can answer.

\section{Discrete random variables}
\labsec{discrete-rv}

\begin{definition}
A \emph{random variable} is a function \(X:\Omega\to\R\). It is \emph{discrete}
when its range is finite or countable, and it is then described by its
\emph{probability mass function} \(p_X(x)=\Pr(X=x)\), which is non-negative and
sums to one.
\end{definition}

Three that recur:

\begin{description}
	\item[Bernoulli\((p\))] \(X\in\{0,1\}\) with \(\Pr(X=1)=p\). The model for a
	single binary label.
	\item[Binomial\((n,p\))] the number of successes in \(n\) independent
	Bernoulli trials, \(p_X(k)=\binom{n}{k}p^{k}(1-p)^{n-k}\).
	\item[Poisson\((\lambda\))] \(p_X(k)=e^{-\lambda}\lambda^{k}/k!\), the limit
	of Binomial\((n,\lambda/n)\) as \(n\to\infty\); counts of rare events.
\end{description}

\section{Continuous random variables}
\labsec{continuous-rv}

\begin{definition}
\(X\) is \emph{continuous} with \emph{density} \(f_X\ge0\) when
\[
	\Pr(a\le X\le b)=\int_a^b f_X(x)\,dx,
	\qquad
	\int_{-\infty}^{\infty}f_X(x)\,dx=1 .
\]
\end{definition}

A density is not a probability: \(f_X(x)\) may exceed one, and
\(\Pr(X=x)=0\) for every single point. What has meaning is \(f_X(x)\,dx\), the
probability of a small interval.

The \emph{cumulative distribution function} \(F_X(x)=\Pr(X\le x)\) works for both
kinds, is non-decreasing, and satisfies \(F_X'=f_X\) in the continuous case.

\begin{description}
	\item[Uniform\((a,b\))] constant density \(1/(b-a)\) on \([a,b]\).
	\item[Exponential\((\lambda\))] \(f_X(x)=\lambda e^{-\lambda x}\) for
	\(x\ge0\); memoryless waiting time.
	\item[Gaussian\((\mu,\sigma^{2}\))]
	\(f_X(x)=\dfrac{1}{\sqrt{2\pi\sigma^{2}}}\exp\Bigl(-\dfrac{(x-\mu)^{2}}{2\sigma^{2}}\Bigr)\).
\end{description}

\marginnote{The Gaussian earns its place twice over: \refch{limit-theorems}
explains why sums of many small effects look Gaussian, and \refch{estimation}
shows that assuming Gaussian noise turns maximum likelihood into least squares.}

\section{Joint, marginal, and conditional distributions}
\labsec{joint}

For two random variables the joint density \(f_{X,Y}\) determines everything.
Integrating out one variable gives a \emph{marginal},
\[
	f_X(x)=\int f_{X,Y}(x,y)\,dy,
\]
and dividing gives a \emph{conditional},
\(f_{Y\mid X}(y\mid x)=f_{X,Y}(x,y)/f_X(x)\) wherever \(f_X(x)>0\). Independence
is the statement \(f_{X,Y}=f_Xf_Y\).

\begin{remark}
Supervised learning is the business of modelling \(f_{Y\mid X}\) and nothing
else. The marginal \(f_X\) governs which inputs are seen, and a change in it
without a change in \(f_{Y\mid X}\) is the mildest form of the distribution shift
named in \refch{map-of-learning-problems}.
\end{remark}

\section{Transforming a random variable}
\labsec{transformations}

If \(Y=g(X)\) with \(g\) differentiable and invertible, the density transforms by
the Jacobian factor of \refsec{jacobian-determinant}:
\[
	f_Y(\vect{y})=f_X\bigl(g^{-1}(\vect{y})\bigr)\,
	\bigl|\det J_{g^{-1}}(\vect{y})\bigr| .
\]

\begin{example}
For \(X\sim\)~Gaussian\((0,1)\) and \(Y=\sigma X+\mu\), the inverse is
\((y-\mu)/\sigma\) with derivative \(1/\sigma\), giving the
Gaussian\((\mu,\sigma^{2})\) density. Standardising a feature is exactly this
change of variables run backwards.
\end{example}

\marginnote{Densities need the Jacobian; expectations do not. Computing
\(\E[g(X)]\) requires no such correction, as \refsec{lotus} explains --- a
distinction worth keeping firmly separate.}

\setchapterpreamble[u]{\margintoc}
\chapter{Expectation, Variance, and Covariance}
\labch{expectation}

\sources{Bertsekas and Tsitsiklis, \emph{Introduction to Probability}, 2nd ed.\
\cite{bertsekas2008probability}, Chapters~2--4; MIT 6.041 \cite{ocw6041},
Lectures~5--12.}

Distributions are unwieldy objects. Expectation reduces one to a single number,
variance to a second, and covariance extends the idea to pairs --- producing, in
the end, the covariance matrix already met in \refch{learning-from-data}.

\section{Expectation}
\labsec{expectation-def}

\begin{definition}
The \emph{expectation} of \(X\) is
\[
	\E[X]=\sum_x x\,p_X(x)
	\qquad\text{or}\qquad
	\E[X]=\int x\,f_X(x)\,dx,
\]
according as \(X\) is discrete or continuous, when the sum or integral converges
absolutely.
\end{definition}

\begin{theorem}[Linearity]
For random variables \(X,Y\) on the same space and constants \(a,b\),
\[
	\E[aX+bY]=a\E[X]+b\E[Y].
\]
No independence is required.
\end{theorem}

\marginnote{Linearity without independence is the workhorse of the subject. It is
what makes the mini-batch gradient of \refsec{sgd} unbiased even though the
examples in a batch share a model and a step.}

\section{Expectations of functions}
\labsec{lotus}

The expectation of \(g(X)\) needs no new distribution:
\[
	\E[g(X)]=\sum_x g(x)p_X(x)
	\qquad\text{or}\qquad
	\E[g(X)]=\int g(x)f_X(x)\,dx .
\]
Contrast \refsec{transformations}: obtaining the \emph{density} of \(g(X)\)
requires a Jacobian factor, obtaining its \emph{expectation} does not.

\begin{remark}
Except for affine \(g\), \(\E[g(X)]\neq g(\E[X])\). For convex \(g\) the
inequality has a direction --- Jensen's inequality gives
\(\E[g(X)]\ge g(\E[X])\) --- which is why an average of losses is not the loss of
an average.
\end{remark}

\section{Variance}
\labsec{variance}

\begin{definition}
\(\Var(X)=\E\bigl[(X-\E[X])^{2}\bigr]\), and the standard deviation is its
square root.
\end{definition}

Expanding the square gives the computational form
\(\Var(X)=\E[X^{2}]-\E[X]^{2}\), and scaling behaves quadratically:
\(\Var(aX+b)=a^{2}\Var(X)\). For \emph{independent} \(X\) and \(Y\),
\(\Var(X+Y)=\Var(X)+\Var(Y)\).

\begin{example}
For \(X\sim\)~Bernoulli\((p)\), \(\E[X]=p\) and \(\E[X^{2}]=p\), so
\(\Var(X)=p-p^{2}=p(1-p)\), largest at \(p=\tfrac12\) and zero at the extremes.
\end{example}

\begin{proposition}
If \(X_1,\ldots,X_n\) are i.i.d.\ with mean \(\mu\) and variance
\(\sigma^{2}\), the sample mean \(\bar X_n=\frac1n\sum_iX_i\) satisfies
\[
	\E[\bar X_n]=\mu,
	\qquad
	\Var(\bar X_n)=\frac{\sigma^{2}}{n}.
\]
\end{proposition}

\marginnote{The \(1/n\) --- and therefore the \(1/\sqrt n\) in the standard
deviation --- is the reason quadrupling a batch halves the gradient noise, and
the reason the returns diminish so quickly.}

\section{Covariance and correlation}
\labsec{covariance-def}

\begin{definition}
\(\Cov(X,Y)=\E\bigl[(X-\E[X])(Y-\E[Y])\bigr]\), and the correlation is
\(\rho=\Cov(X,Y)/\sqrt{\Var(X)\Var(Y)}\in[-1,1]\).
\end{definition}

Independence implies zero covariance; the converse is false. A standard
counterexample is \(X\) symmetric about zero with \(Y=X^{2}\): the two are
completely dependent, yet \(\Cov(X,Y)=\E[X^{3}]=0\).

\begin{remark}
Covariance detects linear association only. Since it is exactly what principal
component analysis in \refsec{pca} consumes, that method sees linear structure
and nothing else --- a limitation of the tool, not of the data.
\end{remark}

\section{The covariance matrix}
\labsec{covariance-matrix}

For a random vector \(\vect{X}\in\R^{d}\) with mean \(\vgreek{\mu}\),
\[
	\Sigma=\E\bigl[(\vect{X}-\vgreek{\mu})(\vect{X}-\vgreek{\mu})\tp\bigr]\in\R^{d\times d}.
\]

\begin{proposition}
\(\Sigma\) is symmetric positive semidefinite, and for any fixed
\(\vect{a}\in\R^{d}\),
\[
	\Var(\vect{a}\tp\vect{X})=\vect{a}\tp\Sigma\vect{a}.
\]
\end{proposition}

\begin{proof}
Symmetry is immediate from the definition. For the second claim, expand
\(\Var(\vect{a}\tp\vect{X})
=\E\bigl[(\vect{a}\tp(\vect{X}-\vgreek{\mu}))^{2}\bigr]
=\vect{a}\tp\E[(\vect{X}-\vgreek{\mu})(\vect{X}-\vgreek{\mu})\tp]\vect{a}\).
Since a variance is non-negative, the quadratic form is non-negative for every
\(\vect{a}\), which is positive semidefiniteness.
\end{proof}

\marginnote{So the population covariance matrix and the sample covariance matrix
\(S\) of \refsec{covariance} are the same object, one exact and one estimated.
The spectral theorem applies to both, which is what principal components rely
on.}

\section{Conditional expectation}
\labsec{conditional-expectation}

\(\E[Y\mid X=x]\) is the mean of the conditional distribution, and
\(\E[Y\mid X]\) is the random variable taking that value when \(X=x\). The tower
property states
\[
	\E\bigl[\E[Y\mid X]\bigr]=\E[Y].
\]

\begin{proposition}
Among all functions \(g\), the one minimising \(\E\bigl[(Y-g(X))^{2}\bigr]\) is
\(g(X)=\E[Y\mid X]\).
\end{proposition}

\begin{remark}
This identifies what a regression model trained on squared error is trying to
be: the conditional mean of the target given the input. Everything the model
gets wrong is either approximation error --- the family cannot express that
conditional mean --- or estimation error from finite data.
\end{remark}

\setchapterpreamble[u]{\margintoc}
\chapter{Limit Theorems}
\labch{limit-theorems}

\sources{Bertsekas and Tsitsiklis, \emph{Introduction to Probability}, 2nd ed.\
\cite{bertsekas2008probability}, Chapter~5; MIT 6.041 \cite{ocw6041},
Lectures~19--20.}

Two theorems justify the entire practice of learning from a sample: averages
converge to what they estimate, and the error of the average has a predictable
shape and size.

\section{Two inequalities}
\labsec{inequalities}

\begin{proposition}[Markov]
If \(X\ge0\) and \(a>0\) then \(\Prob(X\ge a)\le\E[X]/a\).
\end{proposition}

\begin{proof}
\(\E[X]\ge\E[X\cdot\mathbf{1}\{X\ge a\}]\ge a\,\Prob(X\ge a)\), using
non-negativity in the first step.
\end{proof}

\begin{proposition}[Chebyshev]
If \(\E[X]=\mu\) and \(\Var(X)=\sigma^{2}\) then for \(a>0\),
\[
	\Prob\bigl(|X-\mu|\ge a\bigr)\le\frac{\sigma^{2}}{a^{2}}.
\]
\end{proposition}

\begin{proof}
Apply Markov to the non-negative variable \((X-\mu)^{2}\) with threshold
\(a^{2}\).
\end{proof}

\marginnote{These bounds are crude --- they use only two moments --- but they
hold for every distribution, which is exactly what is wanted when the
distribution is unknown.}

\section{The law of large numbers}
\labsec{lln}

\begin{theorem}[Weak law]
Let \(X_1,X_2,\ldots\) be i.i.d.\ with mean \(\mu\) and finite variance
\(\sigma^{2}\). Then for every \(\varepsilon>0\),
\[
	\Prob\bigl(|\bar X_n-\mu|\ge\varepsilon\bigr)\longrightarrow0
	\qquad\text{as }n\to\infty .
\]
\end{theorem}

\begin{proof}
By \refsec{variance}, \(\Var(\bar X_n)=\sigma^{2}/n\). Chebyshev gives
\[
	\Prob\bigl(|\bar X_n-\mu|\ge\varepsilon\bigr)\le\frac{\sigma^{2}}{n\varepsilon^{2}},
\]
which tends to zero.
\end{proof}

\begin{remark}
Training loss is an average over the sample, and test loss is an average over a
different sample from the same distribution. The law is what allows either to
stand in for the expected loss --- and its reliance on the i.i.d.\ hypothesis is
what makes distribution shift so damaging.
\end{remark}

\section{The central limit theorem}
\labsec{clt}

\begin{theorem}
With the same hypotheses and \(\sigma^{2}>0\),
\[
	\frac{\bar X_n-\mu}{\sigma/\sqrt n}
	\ \xrightarrow{\ d\ }\ \text{Gaussian}(0,1).
\]
\end{theorem}

The law of large numbers says the error vanishes; the central limit theorem says
it vanishes at rate \(1/\sqrt n\) and that the rescaled error is Gaussian
regardless of the distribution the samples came from.

\marginnote{This is the honest reason Gaussian noise is assumed so often. It is
not that measurements are Gaussian by nature, but that an accumulation of many
small independent contributions is approximately so.}

\begin{example}
Estimating an accuracy \(p\) from \(n\) independent test examples gives
\(\Var(\bar X_n)=p(1-p)/n\). With \(p\approx0.9\) and \(n=1000\), the standard
deviation is about \(0.0095\), so two systems differing by half a percentage
point on such a test set are not distinguishable by that test set alone.
\end{example}

\section{What the theorems do not say}
\labsec{limit-caveats}

Both require independence and identical distribution, and the central limit
theorem additionally requires finite variance. Heavy-tailed quantities can
violate the last condition, and time-ordered data routinely violate the first.
Where the hypotheses fail, the conclusions are not merely weaker --- they can be
false, and the confidence intervals built on them are then decorative.

\setchapterpreamble[u]{\margintoc}
\chapter{Estimation}
\labch{estimation}

\sources{Bertsekas and Tsitsiklis, \emph{Introduction to Probability}, 2nd ed.\
\cite{bertsekas2008probability}, Chapters~8--9; MIT 6.041 \cite{ocw6041},
Lectures~21--25; MIT 18.05 \cite{ocw1805}.}

Estimation is where probability meets the fitted model. This chapter shows that
the two losses used almost universally in supervised learning are not arbitrary
choices but maximum likelihood under two different noise assumptions --- and
that weight decay is a prior.

\section{The problem}
\labsec{estimation-problem}

Data \(x_1,\ldots,x_n\) are assumed drawn from a family of distributions indexed
by an unknown parameter \(\theta\). An \emph{estimator} is any function
\(\hat\theta(x_1,\ldots,x_n)\) of the data. It is a random variable, so it has a
bias \(\E[\hat\theta]-\theta\) and a variance, and its quality is usually
measured by mean squared error, which decomposes as
\[
	\E\bigl[(\hat\theta-\theta)^{2}\bigr]
	=\bigl(\E[\hat\theta]-\theta\bigr)^{2}+\Var(\hat\theta).
\]

\marginnote{Bias and variance trade against each other. The regularisation of
\refsec{regularisation} buys a large reduction in the second term at the cost of
a small increase in the first --- which is why it works.}

\section{Maximum likelihood}
\labsec{mle}

\begin{definition}
The \emph{likelihood} is the joint density of the observed data read as a
function of \(\theta\),
\[
	L(\theta)=\prod_{i=1}^{n}f(x_i\mid\theta),
\]
and the maximum likelihood estimate is the \(\theta\) maximising it, equivalently
maximising \(\log L(\theta)=\sum_i\log f(x_i\mid\theta)\).
\end{definition}

The logarithm is taken because it converts the product into a sum, which
differentiates cleanly and does not underflow.

\begin{example}
For \(X_i\sim\)~Bernoulli\((p)\),
\(\log L(p)=\sum_i\bigl[x_i\log p+(1-x_i)\log(1-p)\bigr]\). Differentiating and
setting to zero gives \(\hat p=\bar x\), the observed frequency.
\end{example}

\begin{example}
For \(X_i\sim\)~Gaussian\((\mu,\sigma^{2})\) with \(\sigma\) known,
\[
	\log L(\mu)=-\frac{1}{2\sigma^{2}}\sum_i(x_i-\mu)^{2}+\text{const},
\]
so \(\hat\mu=\bar x\). Estimating \(\sigma^{2}\) as well gives
\(\hat\sigma^{2}=\frac1n\sum_i(x_i-\bar x)^{2}\), which is biased low; dividing by
\(n-1\) instead removes the bias.
\end{example}

\section{Two losses, one principle}
\labsec{losses-from-likelihood}

Now let the parameter be the weights of a model \(f_{\vect{w}}\) and condition on
inputs.

\paragraph{Gaussian noise gives squared error.}
Assume \(y_i=f_{\vect{w}}(\vect{x}_i)+\varepsilon_i\) with
\(\varepsilon_i\sim\)~Gaussian\((0,\sigma^{2})\), independent. Then
\[
	\log L(\vect{w})
	=-\frac{1}{2\sigma^{2}}\sum_{i=1}^{n}\bigl(y_i-f_{\vect{w}}(\vect{x}_i)\bigr)^{2}
	+\text{const},
\]
so maximising the likelihood is minimising the sum of squared errors. With a
linear model this is precisely the least-squares problem of
\refch{orthogonality}.

\paragraph{Bernoulli labels give cross-entropy.}
Assume \(y_i\in\{0,1\}\) with
\(\Prob(y_i=1)=\hat y_i=f_{\vect{w}}(\vect{x}_i)\). Then
\[
	-\log L(\vect{w})
	=-\sum_{i=1}^{n}\bigl[y_i\log\hat y_i+(1-y_i)\log(1-\hat y_i)\bigr],
\]
which is the cross-entropy loss.

\begin{remark}
A question left open in \refsec{worked-gradient} is settled here. Squared
error applied to a logistic output carries a factor \(\hat y(1-\hat y)\) that vanishes when the unit
saturates; cross-entropy is the loss that the Bernoulli assumption actually
implies, and its gradient with respect to the pre-activation reduces to
\(\hat y-y\), with no vanishing factor. The improvement is a consequence of
matching the loss to the noise model rather than a numerical trick.
\end{remark}

\section{Bayesian estimation}
\labsec{bayesian}

Maximum likelihood treats \(\theta\) as fixed and unknown. The Bayesian
alternative gives it a distribution: a \emph{prior} \(p(\theta)\) before the data
and a \emph{posterior} after, related by \refeq{bayes},
\[
	p(\theta\mid \text{data})\ \propto\ L(\theta)\,p(\theta).
\]
Reporting the maximiser of the posterior is \emph{maximum a posteriori}
estimation.

\begin{proposition}
For a linear model with Gaussian noise of variance \(\sigma^{2}\) and a Gaussian
prior \(\vect{w}\sim\)~Gaussian\((\vect{0},\tau^{2}I)\), the maximum a posteriori
estimate minimises
\[
	\frac{1}{2\sigma^{2}}\lVert X\vect{w}-\vect{y}\rVert^{2}
	+\frac{1}{2\tau^{2}}\lVert\vect{w}\rVert^{2},
\]
which is ridge regression with \(\delta^{2}=\sigma^{2}/\tau^{2}\).
\end{proposition}

\begin{proof}
Take logarithms in the display above. The log-likelihood contributes the first
term up to a constant, and the log-prior of a Gaussian contributes the second.
Maximising the sum is minimising its negative.
\end{proof}

\marginnote{So the \(\delta^{2}I\) added in \refsec{regularisation} for numerical
reasons, the weight decay added in training for empirical reasons, and a Gaussian
prior on the weights are three descriptions of one modification.}

\section{What this part establishes}
\labsec{part3-summary}

Three conclusions carry into Part~IV. Averages over a sample estimate
expectations, with error shrinking like \(1/\sqrt n\). Squared error and
cross-entropy are maximum likelihood under Gaussian and Bernoulli assumptions
respectively. And regularisation is a prior, which means the choice of penalty is
a statement of belief about the parameters rather than a free knob.


\pagelayout{wide}
\addpart{Deep Learning Foundations}
\pagelayout{margin}

\setchapterpreamble[u]{\margintoc}
\chapter{What a Learning Problem Is}
\labch{learning-problem}

\sources{Zhang, Lipton, Li and Smola, \emph{Dive into Deep Learning}
\cite{d2l2023}, Chapter~1.}

Before any architecture is chosen, a learning problem has to be stated. This
chapter fixes the vocabulary and the four moving parts that every supervised
method rearranges, including the ones this project is preparing for.

\section{Vocabulary}
\labsec{vocabulary}

A few words are used with technical precision throughout Part~IV.

\begin{description}
	\item[Machine learning] fitting a program's behaviour from data instead of
	specifying it by hand --- \emph{programming with data}.
	\item[Deep learning] the subset in which the fitted function is a
	composition of many parameterised layers.
	\item[Dataset] the collection of recorded observations.
	\item[Model] the parameterised family of functions being fitted.
	\item[Parameter] a number adjusted during training.
	\item[Learning algorithm] the procedure that adjusts the parameters.
	\item[Training] the act of running that procedure.
\end{description}

\marginnote{Keeping ``model'' for the family and ``parameters'' for the numbers
inside it avoids a persistent ambiguity: a trained model is a family
\emph{plus} a particular setting, and the two are chosen by different means.}

\section{The four components}
\labsec{four-components}

Any learning problem, however it is later dressed up, consists of four things
that must be specified separately:

\begin{description}
	\item[Data] what the method is allowed to learn from;
	\item[Model] the family of transformations available to it;
	\item[Objective function] a number that says how badly the current model is
	doing;
	\item[Algorithm] the procedure that adjusts the model to reduce that number.
\end{description}

\marginnote{Keeping the four apart is worth the pedantry. Most confusing results
come from changing one of them silently --- a new loss, a new sampling scheme
--- while attributing the difference to the model.}

\subsection{Data}
\labsubsec{data}

A dataset is a collection of \emph{examples}, also called samples or data
points. Each example is described by a fixed list of \emph{features}, also
called covariates, and in supervised problems it carries a \emph{label} --- the
special attribute to be predicted, also called the target. The number of
features is the \emph{dimensionality} of the example, and when that number is the
same for every example the data are of fixed length, which is what makes it
possible to write a batch of \(n\) examples as a matrix \(X\in\R^{n\times d}\) and
reuse everything from \refch{learning-from-data}.

\begin{definition}
Examples are \emph{i.i.d.} --- independent and identically distributed ---
when each is drawn from the same distribution, independently of the others. Nearly every
guarantee about generalisation assumes this, and nearly every deployed system
violates it to some degree.
\end{definition}

The formal content of that definition is \refsec{independence}, and the reason it
matters is \refsec{lln}: without it, an average over the sample stops estimating
anything in particular.

\subsection{Model}
\labsubsec{model}

The model is a parameterised map from features to predictions. Fixing the model
family fixes what the method can express; fixing the parameters within that
family fixes what it currently predicts. For a linear model the reachable
predictions are exactly the column space of the data matrix, which is the
representational question of \refch{vectors-and-column-space} arriving in its
applied form.

\subsection{Objective function}
\labsubsec{objective}

The objective --- the loss or cost function --- turns the gap between prediction
and label into a single number, lower being better. For a real-valued target the
usual choice is squared error,
\[
	\ell(\hat y,y)=\tfrac12(\hat y-y)^{2},
\]
and for a categorical target it is cross-entropy. \refsec{losses-from-likelihood}
shows that these two are not conventions but the maximum likelihood objectives
for Gaussian and Bernoulli noise respectively.

\begin{remark}
The quantity one actually cares about is often not differentiable ---
classification accuracy, ranking quality, revenue. What gets optimised is a
tractable surrogate, and the gap between the surrogate and the real goal is a
permanent source of surprise rather than a temporary inconvenience.
\end{remark}

\subsection{Optimisation algorithm}
\labsubsec{algorithm}

The algorithm searches the parameter space for a setting that makes the
objective small. Gradient descent and its variants dominate, for the practical
reason established in \refsec{backprop}: the gradient of a composed
differentiable model can be computed at roughly the cost of evaluating it, no
matter how many parameters it has.

\section{Training, testing, and overfitting}
\labsec{training-testing}

The available data are split. The \emph{training set} is used to fit the
parameters; the \emph{test set} is held out and used only to estimate how the
fitted model behaves on examples it has never seen.

\marginnote{A test set inspected repeatedly stops being a test set. Each look
leaks information about it into the modelling decisions.}

The split exists because loss on the training set is a biased estimate of
quality. A model flexible enough can drive training loss towards zero by
memorising the particular examples it was given, including their noise.
\emph{Overfitting} is the name for the resulting gap: training loss keeps
falling while test loss turns and rises.

\begin{remark}
Overfitting is not a defect of a specific model family. It is what happens
whenever the capacity of the family is large relative to the information in the
data, and it is the reason that model selection, regularisation, and honest
held-out evaluation are part of the method rather than optional extras.
\end{remark}

By \refsec{clt}, a test set of \(n\) examples estimates the true error with a
standard deviation of order \(1/\sqrt n\). Reported differences smaller than that
are not evidence of anything, however carefully the experiment was run.

\section{Why the framing is worth keeping}
\labsec{framing}

The four components are independent choices, and almost every method discussed
in \refch{map-of-learning-problems} is a different setting of them rather than a
different subject. Recognising that keeps the number of things to learn small:
new architectures change the model, new losses change the objective, and new
optimisers change the algorithm, but the shape of the problem stays put.

\setchapterpreamble[u]{\margintoc}
\chapter{A Map of Learning Problems}
\labch{map-of-learning-problems}

\sources{Zhang, Lipton, Li and Smola, \emph{Dive into Deep Learning}
\cite{d2l2023}, Chapter~1.}

With the four components of \refch{learning-problem} in place, the standard
catalogue of problem types can be laid out as variations on them: what the data
contain, what is being predicted, and whether the learner influences what it
sees next.

\section{Supervised learning}
\labsec{supervised}

Every example carries a label, and the goal is to model the conditional
distribution of the label given the features --- the object \(f_{Y\mid X}\) of
\refsec{joint}.

\subsection{Regression}
\labsubsec{regression}

The label is a real number and the natural loss is squared error, so by
\refsec{conditional-expectation} the fitted function approximates
\(\E[Y\mid X]\). Everything in \refch{orthogonality} applies verbatim when the
model is linear.

\subsection{Classification}
\labsubsec{classification}

The label is one of finitely many \emph{categories} or \emph{classes}. Two
classes make it binary classification; more make it multiclass. The output is a
probability vector, and the loss is cross-entropy, which
\refsec{losses-from-likelihood} derives rather than assumes.

\emph{Hierarchical} classification exploits structure among the classes so that
confusing two neighbouring categories is penalised less than confusing two
distant ones. \emph{Multi-label} classification drops the assumption that the
categories are mutually exclusive, and models each label with its own Bernoulli
output.

\marginnote{Accuracy is a poor summary when classes are unbalanced, for the
reason computed in \refsec{bayes}: with a rare positive class, most positive
predictions can be wrong even at very high accuracy.}

\subsection{Search and ranking}
\labsubsec{ranking}

The output is an ordering of a candidate set rather than a single value, and
quality depends on which items appear near the top. Losses here are built to
approximate rank-based measures, which are themselves not differentiable --- an
instance of the surrogate problem noted in \refsubsec{objective}.

\subsection{Recommender systems}
\labsubsec{recommenders}

Predict a user's response to an item from a sparse, partially observed matrix of
past interactions. The classical treatment is a low-rank factorisation of that
matrix, which is \refsec{low-rank} applied to a table of ratings.

\begin{remark}
Recommendation is the clearest case in which the fitted system changes the data
it will later be trained on: what is recommended is what gets seen, and what
gets seen is what gets recorded. Censoring, incentives, and feedback loops are
open research problems, not solved engineering details.
\end{remark}

\subsection{Sequence learning}
\labsubsec{sequences}

Input and output are sequences whose lengths need not match. Tagging and parsing
attach structure to each position; automatic speech recognition maps audio to
text; text-to-speech maps text to audio; machine translation maps a sequence in
one language to another, generally of different length and different word order.

\section{Unsupervised learning}
\labsec{unsupervised}

No labels are supplied, and the goal is to describe the structure of the data
itself.

\begin{description}
	\item[Clustering] group examples so that members of a group resemble each
	other more than they resemble members of other groups.
	\item[Principal component analysis] find the few directions along which the
	data vary most --- exactly \refsec{pca}.
	\item[Causality and probabilistic graphical models] describe which variables
	influence which, rather than merely which co-vary.
	\item[Generative adversarial networks] fit a mechanism that produces samples
	resembling the data, by training a generator against a discriminator.
\end{description}

\marginnote{Only the second of these is fully settled by linear algebra.
Correlation is what a covariance matrix records, and \refsec{covariance-def}
notes that it detects linear association only --- causal structure is not
recoverable from it without further assumptions.}

\section{Interacting with an environment}
\labsec{environment}

The problems above are \emph{offline}: a fixed dataset is collected first, the
model is fitted second, and the model never influences what it is shown. This is
convenient and often unrealistic.

\begin{definition}
\emph{Distribution shift} occurs when the distribution generating the data at
deployment differs from the one that generated the training data.
\end{definition}

Since \refsec{lln} is what licensed the transfer from sample to population, and
it assumes a single fixed distribution, shift is not a small perturbation of the
theory --- it removes the hypothesis the theory rests on.

\section{Reinforcement learning}
\labsec{reinforcement}

Here the learner is an \emph{agent} placed in a loop with an environment. At
each step it receives an \emph{observation}, chooses an \emph{action} according
to its \emph{policy}, and receives a \emph{reward}; the action influences what is
observed next.

Two difficulties distinguish this from supervised learning. The reward may
arrive long after the action responsible for it, so the agent faces the
\emph{credit assignment} problem. And the agent must balance exploiting what it
already believes to be good against exploring alternatives it has not tried.

\begin{description}
	\item[Markov decision process] the standard model, in which the next state
	depends only on the present state and the chosen action.
	\item[Multi-armed bandit] the one-state case, in which the agent chooses
	repeatedly among fixed options with unknown rewards.
	\item[Contextual bandit] a bandit in which an observation is supplied before
	each choice, but actions still do not affect future states.
\end{description}

\marginnote{The bandit cases are where the estimation machinery of
\refch{estimation} applies most directly: each arm's value is a mean to be
estimated, and the exploration problem is about how confident those estimates
are.}

\section{Terminology carried forward}
\labsec{terminology}

A last group of words appears throughout the literature and is worth fixing
before Part~V: \emph{estimation}; \emph{neural network}; \emph{layer};
\emph{kernel method}; \emph{decision tree}; \emph{graphical model}; and
\emph{representation learning} --- the idea that the useful features should be
learned along with the predictor rather than designed in advance. That last term
is the one \refch{column-space-to-layers} takes up.


\pagelayout{wide}
\addpart{Where the Halves Meet}
\pagelayout{margin}

\setchapterpreamble[u]{\margintoc}
\chapter{From Column Space to Linear Layers}
\labch{column-space-to-layers}

\sources{Strang, \emph{Introduction to Linear Algebra}, 6th ed.\ \cite{strang2023ila},
Chapter~10; Zhang, Lipton, Li and Smola, \emph{Dive into Deep Learning}
\cite{d2l2023}, Chapters~1--5.}

The four parts of this book were assembled separately. This closing chapter says
how they meet, and states --- deliberately conservatively --- what remains to be
studied before the research question that gives the book its title can be
approached.

\section{A network is a composition of the objects already met}
\labsec{network-as-composition}

A feed-forward network alternates two kinds of map. The affine part,
\[
	Z=XW+\vect{1}\vect{b}\tp,
\]
is the batched matrix product of \refsec{batching}; the activation part applies a
fixed scalar function elementwise. Nothing in either is new. What the
composition adds is depth, and depth is what makes the resulting function
non-linear.

\begin{remark}
Stacking affine maps with no activation between them achieves nothing. A
product of matrices is a matrix, so the composition collapses to one linear map
with the same column space. The activation is the only reason a second layer is worth
having, which is a statement about \refsec{four-ways} rather than about biology.
\end{remark}

\section{Three questions, three answers already given}
\labsec{three-questions}

\paragraph{What can the model represent?}
For a single linear layer the reachable outputs are exactly \(\Col(W\tp)\)
applied to the inputs, so the rank of the weight matrix bounds what the layer can
express --- the reachability question of \refsubsec{reachability} in applied
clothing. Compressing a trained layer by truncating its singular values, as in
\refsec{low-rank}, is a deliberate reduction of that rank.

\paragraph{How is it fitted?}
By gradient descent, whose direction is justified by \refsec{steepest}, whose
step size is constrained by the curvature in \refeq{taylor2}, whose convergence
rate is governed by the condition number in \refsec{gd-quadratic}, and whose cost
is made acceptable by the association order in \refsec{backprop}.

\paragraph{Why does the fitted model generalise?}
Because training loss is a sample average of an expectation
(\refsec{lln}), because the deviation of that average shrinks like \(1/\sqrt n\)
(\refsec{clt}), and because the loss itself is a likelihood, so that penalising
the weights is placing a prior on them (\refsec{bayesian}).

\section{Where the linear picture runs out}
\labsec{limits}

Every question above was answered by treating the layer as linear and the
activation as an afterthought. That works because the analysis is local: the
Jacobian of \refsec{matrix-chain-rule} is a genuine linear map, and the whole of
Part~I applies to it at each point. It does not describe the function globally,
and it says nothing about which composition of layers should be preferred.

Two limitations are worth naming, because they are what motivates looking
further:

\begin{description}
	\item[The non-linearity is fixed.] In a conventional layer the parameters
	are the entries of \(W\) and \(\vect{b}\); the activation is chosen in
	advance and not learned. All the flexibility sits in the linear part.
	\item[Interpretation is indirect.] The learned object is a large weight
	matrix whose singular vectors carry no meaning supplied by the problem.
	Understanding what a layer has learned is a separate research effort.
\end{description}

\begin{remark}
Kolmogorov--Arnold networks are motivated by the first of these: they
parameterise the univariate functions themselves rather than the linear mixing.
Whether that reparameterisation delivers the advantages claimed for it is the
question this repository exists to study, and it is not settled by anything in
this book. \textbf{[needs verification]} --- the technical content of that
literature has not yet been read, and no claim about it is made here.
\end{remark}

\section{What to read next}
\labsec{next}

The foundations assembled here are the standard ones, and they are sufficient to
begin. The material that has not been covered, and that the next stage of these
notes should add, is:

\begin{enumerate}
	\item tensors and automatic differentiation in practice, continuing the
	notebooks already in this repository;
	\item approximation theory for univariate functions, in particular spline
	bases, which is the mathematics a Kolmogorov--Arnold parameterisation
	rests on;
	\item the Kolmogorov--Arnold representation theorem itself, in its original
	form, before any neural interpretation of it;
	\item the primary literature on these networks, read against that
	background rather than as a substitute for it.
\end{enumerate}

\marginnote{The order matters. Reading the applied literature before the
representation theorem invites accepting an analogy as a proof --- which is the
failure mode these notes are written to avoid.}

Part~I through Part~IV will be revised as that reading proceeds. The notes are
intended to be maintained rather than finished.


%------------------------------------------------------------------------------
% BIBLIOGRAPHY
%------------------------------------------------------------------------------

\pagelayout{wide}
\nocite{*}
\bibliographystyle{plain}
\bibliography{references/references}
\pagelayout{margin}

\end{document}
